\documentclass[aps,prx,twocolumn,superscriptaddress,longbibliography,nofootinbib,floatfix]{revtex4-2}

\usepackage{amsmath,amssymb,amsthm}
\usepackage{graphicx}
\usepackage[colorlinks=true,urlcolor=blue,linkcolor=blue,citecolor=blue]{hyperref}
\IfFileExists{physics.sty}{\usepackage{physics}}{\DeclareMathOperator{\Tr}{Tr}}
\usepackage[dvipsnames]{xcolor}
\usepackage{bm}
\usepackage{soul}
\IfFileExists{bbold.sty}{\usepackage{bbold}}{}
\usepackage{pifont}
\usepackage{mathtools}
\usepackage{tikz}
\usetikzlibrary{tikzmark,arrows.meta,positioning,shapes.geometric,calc}
\usepackage{array}
\IfFileExists{algorithm2e.sty}{\usepackage[ruled,noline]{algorithm2e}}{}
\usepackage[normalem]{ulem}
\usepackage{mdframed}
\usepackage{enumitem}
\usepackage{makecell}
\usepackage{comment}
\usepackage{booktabs}
\usepackage{microtype}

\DeclareMathOperator{\arctanh}{arctanh}
\DeclareMathOperator{\Sgn}{Sgn}
\DeclareMathOperator{\diag}{diag}

% Shared manuscript shorthands.
\newcommand{\ii}{\boldsymbol{\mathrm{i}}}
\newcommand{\mc}[1]{\mathcal{#1}}
\newcommand{\mr}[1]{\mathrm{#1}}
\newcommand{\mf}[1]{\mathfrak{#1}}
\newcommand{\mb}[1]{\mathbb{#1}}
\newcommand{\mbf}[1]{\mathbf{#1}}
\newcommand{\bs}[1]{\boldsymbol{#1}}
\newcommand{\h}[1]{\hat{#1}}
\newcommand{\Id}{\mathbf{1}}
\newcommand{\cC}{\mc{C}}
\newcommand{\cS}{\mc{S}}
\newcommand{\cU}{\mc{U}}
\newcommand{\hcT}{\h{\mc{T}}}
\newcommand{\hcC}{\h{\mc{C}}}
\newcommand{\hcS}{\h{\mc{S}}}
\newcommand{\hcN}{\h{\mc{N}}}
\newcommand{\comp}[1]{\underset{#1}{\bigcirc}}
\newcommand{\bfr}{{\bf r}}
\newcommand{\TS}{|{\mr{TS}} \rangle\!\rangle}
\newcommand{\CI}{|{\mr{CI}} \rangle\!\rangle}
\newcommand{\vac}{|\mr{vac}\rangle\!\rangle}
\newcommand{\fSWAP}{\mr{fSWAP}}
\newcommand{\AZK}[1]{\mc{K}_{\mr{#1}}}
\newcommand{\cmark}{\textcolor{green}{\ding{51}}}
\newcommand{\xmark}{\textcolor{red}{\ding{55}}}
\newcommand{\drafthighlight}[1]{\textcolor{red}{#1}}
\newcommand{\cmjadd}[1]{\drafthighlight{#1}}
\newcommand{\abadd}[1]{\drafthighlight{#1}}
\newcommand{\HPadd}[1]{\drafthighlight{#1}}
\newcommand{\cmjCom}[1]{\textcolor{red}{[CMJ: #1]}}
\newcommand{\abcom}[1]{\textcolor{purple}{[AB: #1]}}
\newcommand{\HPcom}[1]{\textcolor{cyan}{[HP: #1]}}
\setstcolor{magenta}
\newcommand{\cE}{\mathcal{E}}
\newcommand{\cJ}{\mathcal{J}}
\newcommand{\cK}{\mathcal{K}}
\newcommand{\cT}{\mathcal{T}}
\newcommand{\dd}{\mathrm{d}}
\newcommand{\avg}[1]{\overline{#1}}
% Author-only numerical acceptance gates.  These remain in the source next to
% the statements they govern but are intentionally absent from the rendered
% manuscript.
\newcommand{\ResultPlaceholder}[1]{\ignorespaces}
% Short reader-facing slots for values or conclusions that await production
% data.  Unlike \ResultPlaceholder, these preserve the grammar of the draft.
\newcommand{\ResultValue}[1]{\textcolor{BrickRed}{\textbf{[#1]}}}
\newcommand{\placeholderfigure}[2][0.96\linewidth]{%
  \begingroup
  \setlength{\fboxsep}{8pt}%
  \par\noindent\centering
  \fcolorbox{gray!70!black}{gray!5}{%
    \begin{minipage}[c][0.18\textheight][c]{\dimexpr#1-2\fboxsep-2\fboxrule\relax}
      \centering\small\color{gray!75!black}#2
    \end{minipage}}%
  \par\endgroup}
\newcommand{\appendixplaceholderfigure}[3]{%
  \begin{figure*}[t]
    \placeholderfigure[0.96\textwidth]{\textbf{#2}\\[0.5em]#3}
    \caption{#2.  This is a layout placeholder; the production acceptance
    criteria and complete estimator definitions are given in the text.}
    \label{#1}
  \end{figure*}}

\begin{document}

\title{Chiral Boundary Mode Ensembles in Dynamical Chern Insulators}

\author{Asadullah Bhuiyan}
\affiliation{Department of Physics, Cornell University, Ithaca, New York 14853, USA}
\author{Haining Pan}
\affiliation{Department of Physics and Astronomy, Center for Materials Theory, Rutgers University, Piscataway, New Jersey 08854, USA}
\author{Chao-Ming Jian}
\affiliation{Department of Physics, Cornell University, Ithaca, New York 14853, USA}

\date{\today}

\begin{abstract}
Measurements in an adaptive quantum circuit generate both a conditional quantum state and a classical record drawn from its Born probability.  We study the interplay of these two objects in a finite-range free-fermion circuit whose bulk realizes class-A Chern dynamics and whose local controller programs a one-dimensional domain wall.  The conditional wall states form \ResultValue{STATE-CFT RESULT}, while their propagation is probed in three complementary ways: orientation-resolved modular flow, the full Born linear response including record reweighting, and fixed-record flux insertion.  We further extract \ResultValue{RECORD RESULT} from defect contributions to the Born self-information and a charge-resolved replicated transfer spectrum.  Exact occupation projectors and record-conditioned gain or loss render the usual invertible one-body transfer matrix singular.  We resolve this obstruction by differentiating the normalized covariance map, obtaining a semi-invertible tangent cocycle whose boundary-localized low-exponent sector yields \ResultValue{TANGENT--PURIFICATION RESULT}.  The resulting framework distinguishes the universal physics of the physical trajectory ensemble from both a selected non-Hermitian branch and the generally mixed unconditional channel, and makes a programmable topological interface an operational probe of monitored quantum dynamics.
\ResultPlaceholder{Acceptance-dependent central result: report a stable record effective central charge and sector dimensions only after the replica normalization, bulk subtraction, aspect-ratio collapse, and large-size extrapolation pass the stated acceptance tests.}
\end{abstract}

\maketitle

\ResultPlaceholder{The title and record-CFT headline are retained only
if the replica-calibrated self-information and charge-resolved transfer-spectrum gates in the
numerical campaign both pass.}

\section{Introduction}
\label{sec:introduction}

Measurements do more than decohere a quantum system.  When their outcomes are retained, they generate record-conditioned trajectories whose entanglement and information content can undergo sharp transitions even when ordinary averaged observables remain smooth.  The resulting measurement-induced phases have connected quantum information, nonequilibrium statistical mechanics, and many-body dynamics through problems of entanglement, purification, error correction, and information flow~\cite{LiChenFisher2018,LiChenFisher2019,Skinner2019,GullansHuse2020,JianYouVasseur2020,BaoChoiAltman2020,Fisher2023}.  A defining feature of this setting is that the observed record is not incidental: through the Born rule, it supplies the probability measure on the physical ensemble of quantum trajectories.

Topology makes the role of this ensemble especially sharp.  In equilibrium, a mismatch of bulk invariants is exposed at an interface through protected boundary degrees of freedom; for a Chern insulator, the boundary supports a chiral complex fermion~\cite{Hatsugai1993,HasanKane2010}.  Periodic driving extends bulk--boundary correspondence from a Hamiltonian to an evolution law, while monitored circuits admit symmetry-enriched and topologically ordered trajectory phases with no reduction to a single ground-state problem~\cite{Rudner2013,BaoChoiAltman2021,LavasaniAlaviradBarkeshli2021a,LavasaniAlaviradBarkeshli2021b}.  In one spatial dimension, measurement-only spin circuits can further support symmetry-enriched percolation and gapless SPT steady-state ensembles with protected endpoint modes diagnosed through purification and symmetry-flux string observables~\cite{YuYangLiuLinJian2026}.  Viewed this way, a monitored phase is a dynamical phase in a causally generated temporal-disorder ensemble: each outcome selects the map applied at that time, while the Born rule and adaptive feedback endow the sequence with an intrinsic physical measure and make later maps depend on earlier outcomes.  The relevant robustness question is therefore not only whether topology survives spatial disorder, but whether a trajectory-resolved bulk--boundary correspondence persists under this physical temporal disorder.  This raises a direct dynamical question: when two monitored phases meet, what replaces the equilibrium boundary theory in the physical trajectory ensemble?

For Gaussian fermions, nonunitary circuit evolution can be related to single-particle transfer and localization problems in the Altland--Zirnbauer classes, providing a natural language for monitored topology and criticality~\cite{Jian2022,Kells2023}.  Recent work has constructed programmable defects in monitored free-fermion dynamics and identified topological boundary sectors through Lyapunov spectra~\cite{PanShapourianJian2025,Oshima2025,XiaoKawabata2026}.  These developments establish that topology can be assigned to monitored evolution without first replacing it by an equilibrium Hamiltonian.  They also expose several inequivalent notions of an edge: a mode of a selected transfer problem, a Lyapunov sector of a stochastic product, and a critical state carried by typical physical trajectories.  The immediate starting point for the present work is the adaptive Chern circuit of Bhuiyan, Pan, and Jian, which realizes a postselection-free class-A trajectory ensemble in two dimensions, permits a local domain wall to be programmed into the controller, and exhibits anomalously slow purification at that wall~\cite{BhuiyanPanJian2026}.  What remained unresolved was the theory living on the interface: whether the wall ensemble is critical and chiral, how its slow dynamics survives exact projective rank loss, and whether trajectory-resolved measurement statistics carry universal defect data.

Monitored free fermions display a particularly rich---and subtle---critical landscape.  Continuous and projective monitoring can produce long finite-scale regimes with logarithmic or stronger subvolume entanglement, localization transitions, conserved-charge dynamics, and nonlinear-sigma-model descriptions~\cite{CaoTilloyDeLuca2019,AlbertonBuchholdDiehl2021,FavaPiroliSwann2023,FavaPiroliBernard2024,PoboikoPopperlGornyi2023}.  Above one spatial dimension, related theories support genuine measurement-induced transitions, whereas increasingly large studies of short-range one-dimensional models show that an apparent extended critical regime can cross over only at parametrically large scales~\cite{PoboikoGornyiMirlin2024,NiedereggerVovkStarchl2026}.  A logarithmic entropy fit is therefore not sufficient evidence for the wall considered here.  The same boundary sector must appear consistently in entanglement, charge fluctuations, real-space localization, purification, and orientation-sensitive propagation.

Adaptive circuits provide another nearby perspective.  Measurements, refreshed ancillas, and feedforward can prepare and stabilize entangled targets and can, in suitable constructions, accelerate or reduce the depth of preparation of a common wave function or ground-state sector~\cite{LuLessaKimHsieh2022,Roy2020,Puente2024,Smith2024}.  In that setting, the measurement record often labels a correctable byproduct, and classical coordination can be used to make different branches agree on the desired output.  Interactive dynamics can instead support trajectory-entanglement, steering, control, and absorbing transitions, whose conditioned and averaged critical behavior need not coincide~\cite{RavindranathHanYang2023,ODeaMorningstarGopalakrishnan2024,PiroliLiVasseur2023}; related frustration-free protocols emphasize transport-limited convergence~\cite{Doerstel2026}.  Our protocol lies between these settings.  The correction is local and every measurement outcome is retained, but it does not erase the record or collapse all outcomes onto one wave function; repeated feedback stabilizes an ensemble of inequivalent, charge-fluctuating states that share the same long-distance wall physics.  The wall is consequently a steady defect of the stochastic dynamics, rather than a compiled representative of a predetermined conformal ground state.  The finite-support controller admits a finite-color schedule; Sec.~\ref{sec:model} relates its physical depth to the adaptive preparation bound.

Microscopically, the controller is shaped by the real-space obstruction of a Chern band.  Exponentially localized orthonormal Wannier functions are unavailable, whereas localized nonorthogonal overcomplete frames provide suitable measurement modes; strictly local Gaussian tensor-network representatives face related gaplessness and parent-Hamiltonian obstructions~\cite{Brouder2007,DubailRead2015,LiDongLedwithKhalaf2024}.  We use finite-support approximations to these frames, measure their occupations, and correct an undesired outcome by exchanging a fermion with a freshly prepared occupied or empty ancilla.  Nonorthogonal monitored orbitals can themselves generate anomalous information transport~\cite{PoboikoSzyniszewskiTurner2025}; here their finite-support geometry instead implements a topological controller whose wall is tested against matched trivial circuits.  This ideal perfect-correction step is the fresh-ancilla realization of the earlier adaptive circuit, not an additional stochastic correction probability.  The same local update creates a domain wall either by joining distinct controller frames or by terminating one frame at a chosen interface.  Agreement between these constructions is important because it separates universal wall physics from the microscopic manner in which the control rule is patched.  Native-fermion processors provide several of the required primitives, although a full implementation of the present cycle remains a future objective~\cite{GonzalezCuadraBluvsteinKalinowski2023}.

The physical question requires four related objects to be kept distinct: a selected conditional Kraus word, the probability-weighted ensemble of conditional states, the classical distribution of measurement records, and the unconditional state obtained after discarding those records.  Section~\ref{sec:model} defines all four after constructing the microscopic circuit.  Non-Hermitian Hamiltonians and monitored implementations give powerful descriptions of selected or effective branches, but a branch does not determine the probability-weighted ensemble~\cite{ChaduteauLeeSchindler2026,Fleckenstein2022}.  Indeed, Born-rule and forced-measurement free-fermion circuits can realize different universal behavior even within the same single-particle symmetry class~\cite{JianShapourian2023,Leung2025}.  Probabilistic filters can implement ground-state filtering in suitable settings~\cite{Kosugi2022}.  Our focus is instead the complete Born ensemble as the physical dynamical object.

The unconditional channel provides a complementary, but different, description.  Engineered dissipation can prepare topological steady states and boundary structures~\cite{Bardyn2013,Budich2015,YangMoligniniBergholtz2023}, while finite-range quadratic Lindblad dynamics in more than one dimension cannot converge at a nonvanishing size-independent rate to a unique pure Gaussian topological state under the assumptions of the dissipative preparation obstruction~\cite{Goldstein2019}.  Our construction respects this obstruction in two complementary ways.  With their full exponential tails, the exact OW stabilizers are local but not finite range, which is an allowed route to a common pure Chern target.  At finite $n_{\rm shell}$ the controller is strictly finite range, but truncation destroys the common pure Chern dark state: the conditioned records produce inequivalent trajectories and the discarded-record state $\bar\rho_T$ is generally mixed.  A critical wall may additionally carry a relaxation scale that closes with system size.  Thus neither overcompleteness, fresh ancillas, nor adaptive feedback alone is being invoked as a loophole.  The theorem constrains the unconditional state, whereas the normalized conditional trajectories follow a different evolution; different unravelings of the same averaged evolution can retain different trajectory statistics~\cite{PiccittoRossiniRussomanno2024}.  We therefore compare conditional and unconditional observables directly instead of assigning the trajectory CFT to the averaged channel.

The record itself supplies a third level of description.  Quantum-jump trajectories admit large-deviation and thermodynamic ensembles, while continuous-measurement and fermion-counting theories relate observed stochastic currents to tilted full-counting-statistics generators~\cite{GarrahanLesanovsky2010,Landi2024,Starchl2025,KloiberTollingerSieberer2026}.  More directly, Born-weighted record transfer matrices can encode effective central charges and operator dimensions at measurement-induced critical points~\cite{Zabalo2022}.  We ask whether this construction admits a defect version tied to a programmable class-A interface.  A conformal interpretation requires a replicated transfer problem, subtraction of its extensive bulk and boundary terms, and consistent universal sector data from independent finite-size observables.  This construction yields \ResultValue{RECORD-ONLY BOUNDARY DATA} without reconstructing every conditional state.

We first characterize the state ensemble at the programmed wall.  The numerical observables combine chord-length scaling of the wall excess entropy, quantum interval-charge fluctuations, density correlations, and real-space entanglement contours.  The charge variance is evaluated within each trajectory before the Born average, so its logarithmic coefficient probes the $U(1)$ current algebra rather than the sample-to-sample wandering of the mean charge.  Together these observables determine \ResultValue{STATE-CFT DATA} and localize the critical contribution at the interface.  Starting from a maximally mixed state, the same wall produces \ResultValue{PURIFICATION RESULT}, in contrast with the gapped bulk dynamics.  Because state-level conformal data determine $c_L+c_R$ rather than the propagation direction, we test chirality separately through modular charge drift, the full momentum- and frequency-resolved Born response, and fixed-record flux insertion.  The paired wall orientations provide an internal reversal test for all three probes, while wall-off circuits distinguish directed propagation from a schedule-dependent drift.

We then develop a transfer description adapted to exact projectors.  Differentiating the normalized Gaussian update produces a semi-invertible tangent cocycle that remains well defined when an occupation measurement removes rank.  Its spectrum exhibits \ResultValue{TANGENT-GAP RESULT}, its slow subspace is \ResultValue{LOCALIZATION RESULT}, and its low-exponent density is compared directly with the purification law.  In parallel, we use replicated record weights and charge-resolved transfer sectors to extract \ResultValue{RECORD-CFT RESULT}.  Comparing both constructions with the deterministic averaged covariance map isolates which singular behavior belongs to conditioned trajectories and which survives after the record is discarded.

Taken together, the state, tangent, and record descriptions characterize different aspects of the same programmed interface.  The circuit samples a Born-weighted, charge-fluctuating ensemble, while \ResultValue{SHARED CHIRAL-WALL CONCLUSION} determines the universal structure common to its trajectories.  This distinction is also operational.  A quantum simulator samples the ensemble by retaining every physical record, while the classical record offers a potential route to universal information that can be compared directly with state-tomographic observables.  Figure~\ref{fig:protocol} introduces the adaptive instrument and the two domain-wall constructions.  The following sections then establish the wall state, its orientation, the record theory, the tangent mechanism, and the distinction from unconditional evolution.
\ResultPlaceholder{Headline acceptance test: after deriving the replica normalization, the defect-subtracted record self-information must exhibit a common conformal finite-size collapse, with the same sector data obtained from independent record-transfer estimators.}

\begin{figure*}[t]
  \placeholderfigure[0.97\textwidth]{
    FIGURE 1 PLACEHOLDER\\[0.5em]
    (a) Flattened two-band parent $\h{\mc H}_{\rm CI}$, lower/upper targets, and $A/B$ projected OW modes.\\[0.35em]
    (b) Truncate--then--renormalize construction for the four $3\times3$ controller modes; inset: the displaced overlap zero.\\[0.35em]
    (c) Born occupation measurement, target-dependent fresh ancilla, deterministic $\fSWAP$, and ancilla trace/reset.\\[0.35em]
    (d) One complete cycle: the BPJ-ordered $A-,A+,B-,B+$ word at every active center, together with the four ensemble objects.\\[0.35em]
    (e) Explicit-interface and support-terminated walls, their bond locations, and the orientation convention.\\[0.35em]
    (f) Bulk Chern-marker/twist convergence and matched comparison of the two wall constructions.}
  \caption{
  Adaptive Chern controller and programmed domain walls.  (a) The lower band of the flattened parent at $\alpha_{\rm top}=1$ has $\mathcal C=+1$ and is targeted as occupied, whereas the upper band is targeted as empty.  Our mass parameter has the opposite sign from BPJ.  (b) Each projected mode is multiplied by the real-space window and then normalized independently; $n_{\rm shell}=1$ has $3\times3$-cell support.  (c) The occupation outcome is Born random.  Only a mismatch invokes an occupied or empty fresh ancilla and a deterministic $\fSWAP$; the outgoing ancilla is then discarded without conditioning.  (d) The overlapping channels are visited in the BPJ sequence, lower then upper for $A$ followed by $B$, and form a causal ordered word rather than four commuting stabilizer layers.  (e) The interface may be programmed by patching two local controller frames or by terminating the topological frame.  (f) Topology and wall universality are checked independently.  \ResultPlaceholder{Populate the topology panels only after the Chern-marker and twisted-boundary estimators agree within the stated finite-size tolerance.}}
  \label{fig:protocol}
\end{figure*}

\section{Adaptive circuit and programmed domain walls}
\label{sec:model}

We now specify the adaptive circuit.  We follow the notation and construction sequence of BPJ: first the Chern parent and its projected OW modes, then their occupation stabilizers, and finally the measurement-and-feedforward cycle.  The only band-structure convention change is
\(\alpha_{\rm here}=-\alpha_{\rm BPJ}\), which keeps the numerical input
\(\alpha_{\rm top}=1\) associated with a lower band of Chern number \(+1\).
Figure~\ref{fig:protocol} summarizes the construction.

\subsection{Chern parent, OW modes, and stabilizers}

Consider a square lattice of \(N_x\times N_y\) unit cells with two complex
fermions \(\h c_{\bfr,1}\) and \(\h c_{\bfr,2}\) per cell.  We write
\(\h c(\mbf k)=(\h c_1(\mbf k),\h c_2(\mbf k))^{\mathsf T}\) and define
\begin{align}
 \h{\mc H}_{\rm CI}(\alpha)
 &=\sum_{\mbf k}\h c^\dagger(\mbf k)
 \bigl[\vec n_\alpha(\mbf k)\cdot\vec\sigma\bigr]\h c(\mbf k),
 \nonumber\\[-1mm]
 \vec n_\alpha(\mbf k)
 &=\bigl(\sin k_x,\sin k_y,
 \alpha-\cos k_x-\cos k_y\bigr),
 \nonumber\\[-1mm]
 P_\pm(\mbf k;\alpha)
 &=\frac{1}{2}\left[
 \Id\pm\frac{\vec n_\alpha(\mbf k)}
 {|\vec n_\alpha(\mbf k)|}\cdot\vec\sigma\right].
 \label{eq:parent-bloch-projectors}
\end{align}
Here \(P_+\) and \(P_-\) project onto the upper and lower single-particle
bands, respectively.  With
\begin{equation}
 \mathcal C[P]=\frac{1}{2\pi\ii}\int_{\rm BZ}\dd^2k\,
 \Tr\!\left(P[\partial_{k_x}P,\partial_{k_y}P]\right),
 \label{eq:parent-chern-number}
\end{equation}
the occupied lower band has
\begin{equation}
 \mathcal C[P_-]=
 \begin{cases}
  -1,&-2<\alpha<0,\\
  +1,&0<\alpha<2,\\
  0,&|\alpha|>2.
 \end{cases}
 \label{eq:parent-chern-phases}
\end{equation}
Thus \(\alpha_{\rm top}=1\) gives \(\mathcal C=+1\), while
\(\alpha_{\rm triv}=30\) is trivial.  Equations~\eqref{eq:parent-bloch-projectors}
and \eqref{eq:parent-chern-phases} are BPJ's equations after
\(\alpha_{\rm BPJ}\mapsto-\alpha\); this convention is fixed throughout the
state, flux, and chirality analyses.

Following BPJ, choose the cell-local trial spinors
\begin{equation}
 \tau_A=\frac{1}{\sqrt2}\binom{1}{1},
 \qquad
 \tau_B=\frac{1}{\sqrt2}\binom{1}{-1},
 \label{eq:bpj-trial-spinors}
\end{equation}
and project them into either band:
\begin{equation}
 \lvert\widetilde W_{\bfr,\nu,\pm}\rangle
 =P_\pm(\alpha)\lvert\bfr,\tau_\nu\rangle,
 \qquad \nu=A,B.
 \label{eq:projected-wannier-frame}
\end{equation}
The corresponding normalized localized fermion operator is
\begin{align}
 \h\chi^{(\infty)}_{\bfr,\nu,\pm}
 &=\sum_{\bfr',\mu'}
 W^{(\infty)}_{\bfr,\nu,\pm}(\bfr',\mu')
 \h c_{\bfr',\mu'}
 \nonumber\\
 &\propto\frac{1}{\sqrt{N_xN_y}}\sum_{\mbf k}
 e^{\ii\mbf k\cdot\bfr}
 \tau_\nu^\dagger P_\pm(\mbf k;\alpha)\h c(\mbf k),
 \label{eq:ow-annihilation-operator}
\end{align}
where
\(W^{(\infty)}_{\bfr,\nu,\pm}(\bfr',\mu')
=\langle\widetilde W_{\bfr,\nu,\pm}\vert\bfr',\mu'\rangle/
\|\widetilde W_{\bfr,\nu,\pm}\|\).
These OW modes are localized but not mutually orthogonal.  Their unnormalized
states obey the tight-frame identity
\begin{equation}
 \sum_{\bfr,\nu=A,B}
 \lvert\widetilde W_{\bfr,\nu,\pm}\rangle
 \langle\widetilde W_{\bfr,\nu,\pm}\rvert=P_\pm,
 \label{eq:ow-tight-frame}
\end{equation}
so the two trial families together span the corresponding band.

Define the OW occupation operators
\begin{equation}
 \hcN^{(\infty)}_{\bfr,\nu,\pm}
 =\h\chi^{(\infty)\dagger}_{\bfr,\nu,\pm}
  \h\chi^{(\infty)}_{\bfr,\nu,\pm}.
 \label{eq:ow-number-operator}
\end{equation}
The filled-lower-band state \(\CI\) is uniquely fixed by the BPJ stabilizing
conditions
\begin{equation}
 \hcN^{(\infty)}_{\bfr,\nu,-}\CI=\CI,
 \qquad
 \bigl(\Id-\hcN^{(\infty)}_{\bfr,\nu,+}\bigr)\CI=\CI
 \label{eq:ow-stabilizer-target}
\end{equation}
for every \(\bfr\) and \(\nu=A,B\).  The lower-band modes therefore target
occupation one, whereas the upper-band modes target occupation zero.

\subsection{Finite support and programmed walls}

The physical circuit uses compact approximations to these OW modes.  Let
\(d_{N_\mu}(r_\mu,r'_\mu)\) denote the shortest periodic displacement and,
for \(n=n_{\rm shell}\), define the square window
\begin{equation}
 M_{\bfr}^{(n)}(\bfr')=
 \mathbf{1}\!\left[
 |d_{N_x}(x,x')|\le n,\ 
 |d_{N_y}(y,y')|\le n\right].
 \label{eq:controller-square-window}
\end{equation}
The windowed wavefunction and its operator are
\begin{align}
 W^{(n)}_{\bfr,\nu,\pm}(\bfr',\mu')
 &=
 \frac{M_{\bfr}^{(n)}(\bfr')
 \widetilde W_{\bfr,\nu,\pm}(\bfr',\mu')}
 {\|M_{\bfr}^{(n)}\widetilde W_{\bfr,\nu,\pm}\|},
 \nonumber\\
 \h\chi^{(n)}_{\bfr,\nu,\pm}
 &=\sum_{\bfr',\mu'}W^{(n)}_{\bfr,\nu,\pm}(\bfr',\mu')
 \h c_{\bfr',\mu'}.
 \label{eq:finite-windowed-controller}
\end{align}
Thus the mode is first truncated and only then renormalized.  Before a wall
mask is imposed, \(n_{\rm shell}=1\) has exactly \(3\times3\) cells of support,
or at most 18 microscopic orbitals.  At \(\alpha=1\), it retains \(0.992155\)
of the untruncated squared norm and a target-band overlap zero with winding
\(\nu_{\rm ov}=-1\), displaced from \(\mbf k/\pi=(0.5,0)\) to approximately
\((0.4922,0)\).  In our convention
\(\mathcal C[P_-]=-\nu_{\rm ov}=+1\).  This is an overlap-obstruction
diagnostic, not a Chern number of the compact spinor, which defines a smooth
global section and has \(\mathcal C_{\rm compact}=0\).
Appendix~\ref{app:finite-support-controller} gives the convolution and winding
analysis.  Because windowing introduces out-of-band weight, the compact
occupations no longer share the exact common eigenstate in
Eq.~\eqref{eq:ow-stabilizer-target}; the finite-range circuit stabilizes an
ensemble rather than one exact parent-Hamiltonian ground state.
Consequently, the finite-shell cutoff is an approximation at the level of
the local controller, not a compactly supported representative of the Chern
bundle.  Its large retained norm and nonzero overlap winding preserve the
imprint of the target projector, but not the exactly flat topological dark
subspace excluded for a finite-range dissipative preparation with a unique
pure steady state and nonvanishing thermodynamic relaxation rate
\cite{Goldstein2019}.

We program the wall in two ways.  Let
\(\Theta_{\rm top}(\bfr)=\mathbf{1}[x_L\le x\le x_R]\) and
\begin{equation}
 \alpha_{\bfr}=
 \begin{cases}
  \alpha_{\rm top},&\Theta_{\rm top}(\bfr)=1,\\
  \alpha_{\rm triv},&\Theta_{\rm top}(\bfr)=0.
 \end{cases}
 \label{eq:controller-slab}
\end{equation}
For the \emph{explicit interface}, every center remains active and
Eq.~\eqref{eq:finite-windowed-controller} is formed using the uniform-parent
projector \(P_\pm(\mbf k;\alpha_{\bfr})\) assigned to that center.  Compact
tails may cross the interface.  This construction patches locally uniform
controller frames; it is not the spectral projector of one inhomogeneous
Hamiltonian.

For the \emph{support-terminated wall}, define
\begin{equation}
 D_{\bfr}(\bfr')=
 \Theta_{\rm top}(\bfr)\Theta_{\rm top}(\bfr')
 +[1-\Theta_{\rm top}(\bfr)][1-\Theta_{\rm top}(\bfr')]
 \label{eq:same-region-mask}
\end{equation}
and, for active centers \(\bfr\) in the topological slab, use
\begin{equation}
 W^{(n,{\rm term})}_{\bfr,\nu,\pm}
 =
 \frac{D_{\bfr}M_{\bfr}^{(n)}
 \widetilde W_{\bfr,\nu,\pm}(\alpha_{\rm top})}
 {\|D_{\bfr}M_{\bfr}^{(n)}
 \widetilde W_{\bfr,\nu,\pm}(\alpha_{\rm top})\|}.
 \label{eq:support-terminated-wall}
\end{equation}
The diagonal masks commute and are both imposed before one normalization; the
order written in Eq.~\eqref{eq:support-terminated-wall} has no physical
content.  Only topological-slab centers are subsequently updated.  The
exterior is prepared as a Born-sampled onsite product state and is excluded
from the OW word.  We denote its stored preparation record by \(\zeta\), its
probability by \(p_\zeta\), and the resulting initial state by
\(\rho_{0|\zeta}\).  For the explicit-interface circuit, \(\zeta\) is empty
and \(p_\zeta=1\).  Periodicity in \(x\) places oppositely oriented walls on
the bonds \(x_L-\tfrac12\) and \(x_R+\tfrac12\).  The globally diagonalized
inhomogeneous exact-\(H\) problem is a third, deterministic benchmark and is
neither circuit construction.  We find
\ResultValue{AGREEMENT OF THE TWO WALL CONSTRUCTIONS} after transverse
convergence [Fig.~\ref{fig:protocol}(e,f)].

\subsection{Perfect correction and one complete BPJ-ordered cycle}

We now define the stochastic update.  For a normalized controller mode
\(\h\chi_j=\sum_i\chi_{j,i}c_i\), where
\(j=(\bfr,\nu,\sigma)\) and \(\sigma=\pm\), let
\(\hcN_j=\h\chi_j^\dagger\h\chi_j\),
\(\Pi_{j,1}=\hcN_j\), and \(\Pi_{j,0}=\Id-\hcN_j\).  A Gaussian state is
represented by
\begin{align}
 C_{ij}&=\Tr(\rho c_i^\dagger c_j),
 &G&=2C-\Id,
 \nonumber\\
 p_{j,m}(G)&=\frac{1+\eta_m\chi_j^\dagger G\chi_j}{2},
 &\eta_m&=2m-1.
 \label{eq:covariance-and-born-rule}
\end{align}
The target occupations are \(s_-=1\) and \(s_+=0\).  Conditional on the Born
outcome \(m\), perfect correction uses
\begin{equation}
 \begin{array}{c|cc}
  &m=0&m=1\\ \hline
  s=0&K_{0|0}=\Pi_0&K_{1|0}=\h\chi\\[1mm]
  s=1&K_{0|1}=\h\chi^\dagger&K_{1|1}=\Pi_1
 \end{array},
 \qquad
 \rho'_{m|s}=\frac{K_{m|s}\rho K_{m|s}^\dagger}{p_m}.
 \label{eq:main-perfect-correction-kraus}
\end{equation}
The channel index \(j\) is suppressed in the table.  The identities
\(K_{m|s}^\dagger K_{m|s}=\Pi_m\),
\(\Pi_sK_{m|s}=K_{m|s}\), and
\(\sum_mK_{m|s}^\dagger K_{m|s}=\Id\) show that the correction preserves the
Born probability and resets the measured OW mode to its target occupation.

This ideal model is obtained from a local, number-conserving dilation of the
BPJ correction step.  On a mismatch, prepare a fresh ancillary mode
\(\h d\) in \(\lvert s\rangle_d\), apply
\begin{align}
 \fSWAP(\h\chi,\h d)
 &=\exp\!\left(\ii\frac{\pi}{2}\h X^\dagger\h X\right),
 &\h X&=\h\chi-\h d,
 \nonumber\\
 \lvert m\rangle_\chi\lvert s\rangle_d
 &\longmapsto
 \lvert s\rangle_\chi\lvert m\rangle_d.
 \label{eq:main-fswap-dilation}
\end{align}
and discard or reset the outgoing factorized ancilla.  No favorable outcome
is selected.  Total charge is conserved on the enlarged physical--ancilla
system even though the reduced physical trajectory contains gain and loss.
Unlike BPJ's reusable ancillary layer, the idealization integrates out a
fresh target-state ancilla after every mismatch and therefore has no separate
bath-redistribution stage.

For later use, write \(P_j=\chi_j\chi_j^\dagger\),
\(Q_j=\Id-P_j\).  The exact conditional covariance map is
\begin{equation}
 \boxed{\Phi^{(j)}_{m|s}(G)=
 Q_j\bigl(\Id+\eta_mGP_j\bigr)^{-1}GQ_j+\eta_sP_j.}
 \label{eq:conditional-covariance-main}
\end{equation}
The inverse is a rank-one resolvent and exists on every nonzero-probability
branch.  Equation~\eqref{eq:conditional-covariance-main}, rather than an
unnormalized non-Hermitian propagator, is differentiated in
Sec.~\ref{sec:tangent}.

\begin{mdframed}[nobreak=true,linewidth=0.5pt,innerleftmargin=5pt,
 innerrightmargin=5pt,innertopmargin=4pt,innerbottommargin=4pt]
\noindent\textbf{Algorithm 1: One perfect-correction cycle.}
\begin{enumerate}[leftmargin=1.35em,itemsep=1pt,topsep=3pt]
 \item Choose the ordered active centers
 \(\pi_t=(\bfr_{t,1},\ldots,\bfr_{t,N_{\mathcal A}})\).
 \item For each center, visit \(\nu=A,B\).  First measure
 \(\hcN_{\bfr,\nu,-}\); if \(m_-=0\), supply an occupied fresh ancilla and
 apply \(\fSWAP\).
 \item Next measure \(\hcN_{\bfr,\nu,+}\); if \(m_+=1\), supply an empty fresh
 ancilla and apply \(\fSWAP\).
 \item Update \(G\) with Eq.~\eqref{eq:conditional-covariance-main} after
 every outcome, and store the complete causal record together with \(\pi_t\).
\end{enumerate}
\end{mdframed}

Thus the within-center word is exactly the BPJ sequence
\begin{align}
 \mathcal Q_{\rm BPJ}
 &=\bigl((A,-),(A,+),(B,-),(B,+)\bigr),
 \nonumber\\
 (s_q)_{q\in\mathcal Q_{\rm BPJ}}&=(1,0,1,0).
 \label{eq:bpj-channel-word}
\end{align}
Let \(j=1,\ldots,D_T\) enumerate the resulting global chronological word.
For prescribed center words
\(\pi=(\pi_1,\ldots,\pi_T)\), exterior record \(\zeta\), and OW outcomes
\(\xi\), the complete operator and causal probability are
\begin{align}
 K_{\xi|\pi}
 &=\overleftarrow{\prod_{t=1}^{T}}
   \overleftarrow{\prod_{\ell=1}^{N_{\mathcal A}}}
   \overleftarrow{\prod_{q\in\mathcal Q_{\rm BPJ}}}
   K^{(\bfr_{t,\ell},q)}_{m_{t\ell q}|s_q},
 \nonumber\\
 p_{\xi|\pi,\zeta}
 &=\Tr(K_{\xi|\pi}\rho_{0|\zeta}K_{\xi|\pi}^\dagger)
 =\prod_{j=1}^{D_T}
 \frac{1+\eta_{m_j}\chi_j^\dagger G_{<j|\zeta}\chi_j}{2}.
 \label{eq:causal-cycle-word}
\end{align}
The arrows mean that later operations act on the left, and each Born factor is
evaluated on the covariance produced by all preceding operations.

The four ensemble objects introduced in Sec.~\ref{sec:introduction} are
\begin{equation}
 \begin{aligned}
  \rho_{\zeta,\xi|\pi}
  &=\frac{K_{\xi|\pi}\rho_{0|\zeta}K_{\xi|\pi}^\dagger}
  {p_{\xi|\pi,\zeta}},\\
  \bar\rho_{T|\pi}
  &=\sum_{\zeta,\xi}p_\zeta p_{\xi|\pi,\zeta}
  \rho_{\zeta,\xi|\pi},\\
  \mathbb P(\pi,\zeta,\xi)
  &=P_{\rm sch}(\pi)p_\zeta p_{\xi|\pi,\zeta},
  &\bar\rho_T&=\sum_\pi P_{\rm sch}(\pi)\bar\rho_{T|\pi}.
 \end{aligned}
 \label{eq:record-K-p}
\end{equation}
Thus \(K_{\xi|\pi}\) is a selected conditional OW word;
\(\{p_\zeta p_{\xi|\pi,\zeta},\rho_{\zeta,\xi|\pi}\}\) is the physical Born
ensemble at fixed schedule; \(p_\zeta p_{\xi|\pi,\zeta}\) is its classical
record law; and \(\bar\rho_{T|\pi}\) is the discarded-record state.  The
physical outcomes \((\zeta,\xi)\) constitute causal temporal disorder
conditional on \(\pi\).  If the center ordering is randomized,
\(P_{\rm sch}\) describes additional externally chosen disorder and must be
stored separately.  Because neighboring OW projectors overlap, their order
is physical: randomized center words are the production choice, while raster
and reverse-raster center words are controls.

The controller modes are derived from flattened uniform-parent projectors,
but truncation, chronological measurement, Born sampling, and feedback do
not equal evolution under the exact domain-wall Hamiltonian.  Dashed
exact-\(H\) curves use its ground-state covariance only as a matched finite
geometry benchmark.  Real-space markers and twisted-boundary estimators give
\ResultValue{CHERN/BOTT CONVERGENCE AND $n_{\rm shell}=1$ ERROR}
[Fig.~\ref{fig:protocol}(f)]~\cite{BiancoResta2011,BauMarrazzo2024,
FukuiHatsugaiSuzuki2005}.  Appendix~\ref{app:adaptive-depth-bound} gives the
finite-color compilation and the adaptive depth bound.  For stationary
observables, every production trajectory undergoes its own size-dependent
convergence interval before its observation window is collected, and
already Born-sampled trajectories enter ensemble averages with equal Monte
Carlo weight.

\section{Critical state ensemble at the wall}
\label{sec:state-cft}

We first characterize the steady-state wall independently of its direction of propagation.  Figure~\ref{fig:state-cft} combines interval entanglement, $U(1)$ charge fluctuations, real-space contours, correlations, and purification to test whether the Born trajectories share a one-dimensional critical boundary theory.

For a wall interval of length $\ell$, we use the chord distance
\begin{equation}
 d_L(\ell)=\frac{L}{\pi}\sin\!\left(\frac{\pi\ell}{L}\right),
 \label{eq:chord}
\end{equation}
and define the wall-excess entropy by subtracting a matched bulk or wall-off geometry before fitting
\begin{equation}
 \Delta S(\ell)=a_S\log d_L(\ell)+b_S+\delta_S(\ell,L).
 \label{eq:entropy-fit}
\end{equation}
The conversion between $a_S$ and the central charge is fixed by whether the chosen cut contains one wall branch, two oppositely oriented walls, or a boundary interval.  We use that geometry without changing the convention between the exact-$H$ benchmark and the Born ensemble.  Figure~\ref{fig:state-cft}(a) gives \ResultValue{$a_S$, $c_{\rm edge}$, AND LARGE-$L$ DRIFT}.  This entropy scaling determines the total critical content seen by the cut, $c_L+c_R$, but not its orientation~\cite{ChenLiFisher2020,LiChenLudwig2021}.

Charge cumulants supply independent state data for free fermions~\cite{KlichLevitov2009,SongFlindtRachel2011}.  For the $U(1)$ boundary theory, the current-algebra normalization derived in Appendix~\ref{app:estimators} relates the logarithmic interval variance to its current level.  With $N_A=\sum_{j\in A}n_j$, we first form the quantum variance inside each trajectory and only then average over trajectory-resolved records,
\begin{align}
 F_{A,\xi}&=\big\langle N_A^2\big\rangle_\xi-\big\langle N_A\big\rangle_\xi^2,\nonumber
 \\
 \overline{F_A}&=a_Q\log d_L(\ell)+b_Q+\delta_Q(\ell,L)
\label{eq:charge-fit}
\end{align}
which probes the $U(1)$ current algebra rather than repeating the entropy fit.  Figures~\ref{fig:state-cft}(b) and \ref{fig:state-cft}(c) yield \ResultValue{$a_Q$/CURRENT LEVEL, CORRELATION EXPONENT, AND WALL CONTOUR WEIGHT}.  The record-to-record variance of $\langle N_A\rangle_\xi$ is a different ensemble observable, familiar from charge-sharpening problems, and is not substituted for the within-trajectory quantum variance~\cite{AgrawalZabaloChen2022,BarrattAgrawalGopalakrishnan2022}.  The entanglement and charge contours then verify in real space that the logarithmic excess is carried by the programmed walls.  Along the wall, we report density and current correlations through their decay exponents; a current correlator proportional to $r^{-2}$ has scaling dimension $\Delta=1$.

We next probe how this state is reached.  Starting from the maximally mixed covariance $C_0=\Id/2$, we monitor the entropy or purity deficit $\Delta S_{\mathrm{pur}}(t)$ for wall-on and wall-off circuits.  Figures~\ref{fig:state-cft}(d) and \ref{fig:state-cft}(e) show \ResultValue{WALL POWER LAW, BULK RATE, AND TRANSVERSE LOCALIZATION}.  The tangent spectrum in Sec.~\ref{sec:tangent} relates this slow decay to \ResultValue{MATCHED BOUNDARY TANGENT SECTOR}.  \ResultPlaceholder{A claimed power law must survive changes of fit onset, fit window, aspect ratio, and initial covariance.}

\begin{figure*}[t]
  \placeholderfigure[0.97\textwidth]{
    FIGURE 2 PLACEHOLDER\\[0.5em]
    (a) Wall-excess entropy versus chord length, with exact-H dashed benchmark and fit-window audit.\\[0.35em]
    (b) Connected charge fluctuation and the independently extracted $U(1)$ coefficient; inset: wall correlator and decay exponent.\\[0.35em]
    (c) Entanglement/charge contour showing localization of the logarithmic contribution at the two programmed walls.\\[0.35em]
    (d) Wall-on versus wall-off purification over time and system size, including effective exponents.\\[0.35em]
    (e) Transverse profile of the residual late-time entropy or purity defect, identifying the wall as the slow sector.}
  \caption{
  Critical state ensemble and slow wall purification.  (a) Wall-excess entropy versus chord length, compared with the exact-$H$ result on the same finite geometry.  (b) Within-trajectory interval charge variance and its $U(1)$ current coefficient, with the wall correlation decay shown in the inset.  (c) Entanglement and charge contours resolving the two programmed walls.  (d) Purification from the maximally mixed state for wall-on and wall-off circuits.  (e) Transverse profile of the late-time mixedness.  Trajectory estimators are averaged with their Born-sampled weights, and dashed curves are benchmarks rather than fit constraints.  Entanglement fixes $c_L+c_R$ for this geometry, not the sign of the chiral central charge.  \ResultPlaceholder{Accept the state-CFT claim only if the entropy and charge coefficients extrapolate consistently, the contour remains wall localized, both wall constructions agree, and the inferred coefficient does not drift downward at the largest sizes.}  \ResultPlaceholder{Accept algebraic purification only if a tangent-spectrum prediction from Fig.~\ref{fig:tangent} accounts for the same late-time exponent.}}
  \label{fig:state-cft}
\end{figure*}

\section{Chiral dynamics of the domain wall}
\label{sec:chirality}

Having established the critical content of the wall, we next determine its direction of propagation.  A $c=1$ logarithm alone does not fix $c_-=c_R-c_L$, so Fig.~\ref{fig:chirality} brings together three complementary, orientation-sensitive probes.

\subsection{Modular charge drift}

For a trajectory covariance restricted to a wall region $A$, let $K_A=-\log\rho_A$ be the Gaussian modular Hamiltonian.  A localized charge perturbation is evolved in modular time $s$,
\begin{equation}
 \delta n_\xi(y,s)=
 \Tr\!\left[\delta\rho_{A,\xi}\,
 e^{\ii sK_{A,\xi}}n_y e^{-\ii sK_{A,\xi}}\right].
 \label{eq:modular-evolution}
\end{equation}
We extract its centroid only while the packet is localized and before it meets the second wall or subsystem endpoint.  We require the drift to reverse between oppositely oriented walls and to vanish in wall-off and exact-trivial controls.
Figure~\ref{fig:chirality}(a) gives \ResultValue{SIGNED MODULAR VELOCITIES AND WALL-OFF NULL RESULT}.  Because modular time is generated by the reduced state rather than the physical circuit, this drift measures the handedness of the state but is not a quantized transport coefficient.

\subsection{Full Born response, including the score}

The physical response of the Born ensemble includes both the change of the conditional state and the change in the probability of obtaining its record.  For a localized potential impulse $V$ and a trajectory observable $O_\xi[V]$,
\begin{align}
 \left.\frac{\delta\langle O\rangle_V}{\delta V}\right|_{V=0}
 &=\left\langle
 \left.\frac{\delta O_\xi[V]}{\delta V}\right|_{V=0}
 \right\rangle_\xi
 +\left\langle \delta O_\xi\,s_\xi\right\rangle_\xi,
 \label{eq:born-response}\\
 s_\xi&=\left.\frac{\delta\log p_\xi[V]}{\delta V}\right|_{V=0},
 \qquad
 \delta O_\xi=O_\xi-\langle O\rangle.
 \label{eq:score}
\end{align}
The second term is the likelihood-score contribution.  Omitting it gives the averaged conditional-state response, rather than the response of the physical Born ensemble.  With $O=n(y,t)$, we Fourier transform along the wall and in circuit time.  Figure~\ref{fig:chirality}(b) shows \ResultValue{ONE-WAY RIDGE, $v$, ORIENTATION REVERSAL, AND SCORE CONTRIBUTION}.  A low-energy feature $\omega\simeq vk$ whose slope reverses with the wall orientation is the direct monitored-time propagation test in our analysis.

\subsection{Fixed-record twist flow}

Flux insertion supplies a complementary branchwise diagnostic.  We freeze one realized circuit record $\xi$, insert a $U(1)$ twist $\phi$ through the periodic direction of that same realization, and follow the wall entanglement levels
\begin{equation}
 \varepsilon_a(\phi)=\log\!\frac{1-\nu_a(\phi)}{\nu_a(\phi)}
 \label{eq:twist-entanglement-levels}
\end{equation}
continuously for $0\leq\phi\leq2\pi$.  Here $\nu_a(\phi)$ are the overlap-tracked eigenvalues of the wall-restricted trajectory covariance.  The twist multiplies only controller modes crossing a fixed seam by $e^{\ii\phi}$; the stochastic couplings, update order, and record are otherwise unchanged.  Figure~\ref{fig:chirality}(c) yields \ResultValue{ORIENTED SPECTRAL FLOW AND TRIVIAL CONTROL}.  Independently resampling the record at every $\phi$ would not define a continuous family and would have no Berry-holonomy interpretation.  Fixed-record flow is therefore an auxiliary branchwise topology diagnostic, rather than the Born response of Eqs.~\eqref{eq:born-response} and \eqref{eq:score}; realizing it experimentally would require correlated record control.

\begin{figure*}[t]
  \placeholderfigure[0.97\textwidth]{
    FIGURE 3 PLACEHOLDER --- THREE COMPLEMENTARY CHIRALITY PROBES\\[0.6em]
    (a) Modular charge-packet trajectories for the two wall orientations, with wall-off control.\\[0.45em]
    (b) Full Born response $\chi_{\mathrm{Born}}(k,\omega)$, decomposed into conditional and likelihood-score terms; oriented low-energy ridge.\\[0.45em]
    (c) Fixed-record $0\!\to\!2\pi$ twist flow for paired wall orientations and matched trivial control.}
  \caption{
  Three complementary tests of chirality.  (a) Modular charge-packet motion in the two wall orientations and the wall-off control.  (b) Momentum- and frequency-resolved Born susceptibility, separated into conditional-state and likelihood-score contributions.  (c) Fixed-record twist flow for the paired wall orientations and a matched trivial circuit.  Modular drift diagnoses directionality in the reduced state, the Born susceptibility measures physical monitored-time propagation, and the twist flow probes the topology of a continuous branch family.  \ResultPlaceholder{Use ``chiral'' in the title only if at least two orientation-sensitive diagnostics survive size scaling and reversal tests, with the full Born response providing a one-way low-energy feature or the twist flow providing robust oriented spectral flow.}}
  \label{fig:chirality}
\end{figure*}

\section{Universal data in trajectory-resolved records}
\label{sec:record-cft}

We next ask whether the same wall theory can be identified from the classical measurement record alone.  The relation
\begin{equation}
 Z_\xi\equiv p_\xi
 \label{eq:record-partition}
\end{equation}
is only a starting point and does not by itself imply conformal structure.  The directly sampled observable is the record self-information
\begin{equation}
 I_\xi=-\log p_\xi,
 \qquad H_{\rm rec}=\langle I_\xi\rangle_\xi,
 \label{eq:record-self-information}
\end{equation}
which is a Shannon entropy.  To connect it to a statistical-mechanics free energy, we introduce the replica family
\begin{equation}
 \mathcal Z_n(L,T)=\sum_\xi p_\xi^n,
 \qquad
 -\left.\partial_n\log\mathcal Z_n\right|_{n=1}=H_{\rm rec}.
 \label{eq:record-replica}
\end{equation}
We derive the transfer representation and normalization of this replica derivative before assigning a conformal interpretation.  The matched combination $H_{\rm rec}^{\mathrm{def}}$ uses wall-on, wall-off, and homogeneous circuits with identical local update counts.  This subtraction cancels the common spacetime bulk contribution, while leaving a nonuniversal wall free-energy density that must be fitted together with the subextensive term.

Let $r=T/L$ be the temporal aspect ratio.  Dividing the matched combination by the number of walls, the wall length, and the observation time defines
\begin{equation}
 f_{\rm def}(L,r)
 =\frac{H_{\rm rec}^{\rm def}(L,rL)}{N_{\rm wall}L(rL)}
 =f_\infty(r)-\frac{A_{\rm rec}(r)}{L^2}
 +\frac{b_4(r)}{L^4}+\cdots .
 \label{eq:record-casimir}
\end{equation}
Equation~\eqref{eq:record-casimir} is the agnostic fit performed before any CFT conversion.  Only after the replica normalization fixes a chiral periodic cylinder and the large-$r$ limit is stable do we identify the per-wall amplitude as
\begin{equation}
 A_{\rm rec}(r)\xrightarrow{r\to\infty}
 \frac{\pi v_{\rm rec}c_{\rm eff}^{\rm rec}}{12}.
 \label{eq:record-amplitude}
\end{equation}
The numerical coefficient is rederived if the implemented spacetime boundary conditions differ from this cylinder.  We determine the nonuniversal velocity $v_{\rm rec}$ independently from transfer gaps or aspect-ratio collapse.  Figures~\ref{fig:record-cft}(a) and \ref{fig:record-cft}(b) give \ResultValue{$f_\infty$, $A_{\rm rec}$, INDEPENDENT $v_{\rm rec}$, AND $c_{\rm eff}^{\rm rec}$}.  This effective central charge belongs to a replicated, generally nonunitary record theory and is not set equal to the trajectory-state value inferred from Fig.~\ref{fig:state-cft}.

Charge-resolved record transfer eigenvalues provide a second finite-size observable.  If $\Lambda_a(L)$ labels a sector of the record transfer problem, the expected form is
\begin{equation}
 -\log\left|\frac{\Lambda_a(L)}{\Lambda_0(L)}\right|
 =\frac{2\pi v_{\mathrm{rec}}}{L}x_a^{\mathrm{rec}}+o(L^{-1}).
 \label{eq:record-gaps}
\end{equation}
Figure~\ref{fig:record-cft}(c) yields \ResultValue{CHARGE SECTORS, COMMON VELOCITY, AND $x_a^{\rm rec}$}.  Agreement of $v_{\mathrm{rec}}$ between Eqs.~\eqref{eq:record-amplitude} and \eqref{eq:record-gaps} prevents an arbitrary rescaling from manufacturing a central charge.  The sector labels are defined from conserved charge or a controlled record deformation, rather than by sorting noisy eigenvalues after the fact.

This construction builds on trajectory thermodynamics and the trajectory-resolved statistical mechanics of monitored circuits, where record entropies and replicated transfer gaps can encode universal information~\cite{GarrahanLesanovsky2010,JianYouVasseur2020,BaoChoiAltman2020,Zabalo2022}.  Here the transfer problem carries a defect tied to a programmable class-A interface.  Related cross-entropy and postselection-free measurement strategies motivate operational estimators that avoid a histogram over exponentially many complete records~\cite{LiZou2023,GarrattAltman2024}.  Figure~\ref{fig:record-cft}(d) reports \ResultValue{FINITE-SAMPLE ESTIMATOR BIAS, ERROR, AND SAMPLE SCALING} against the exactly evaluated simulated likelihoods.  \ResultPlaceholder{The operational record-only claim is accepted only after the replica mapping and finite-sample estimator are both numerically validated; direct access to simulated $p_\xi$ alone is insufficient.}

\begin{figure*}[t]
  \placeholderfigure[0.97\textwidth]{
    FIGURE 4 PLACEHOLDER\\[0.5em]
    (a) Matched defect self-information density for several $L$ and $T/L$, with $f_\infty$ and the $L^{-2}$ term shown separately.\\[0.35em]
    (b) Casimir/aspect-ratio collapse and extraction of $v_{\mathrm{rec}}c_{\mathrm{eff}}^{\mathrm{rec}}$.\\[0.35em]
    (c) Charge-resolved record transfer gaps and common-velocity extraction of $x_a^{\mathrm{rec}}$.\\[0.35em]
    (d) Record-only finite-sample estimator benchmark against exact simulated likelihoods.}
  \caption{
  Universal data in trajectory-resolved records.  (a) Matched defect self-information density and its $L^{-2}$ coefficient for several temporal aspect ratios.  (b) Cylinder collapse and the independent extraction of the record velocity and effective central charge.  (c) Charge-resolved record transfer gaps and sector dimensions.  (d) Finite-sample record-only estimator benchmarked against exact simulated likelihoods.  The state central charge and the record effective central charge are displayed separately throughout.  \ResultPlaceholder{An extensive fit with a small residual is not sufficient: the residual must have the predicted aspect-ratio dependence, the independently determined velocity, stable sector assignments, and a wall-off null control. Insert the record-CFT headline values only if all four panels pass; otherwise retain Fig.~\ref{fig:record-cft} as a null or crossover result.}}
  \label{fig:record-cft}
\end{figure*}

\section{Tangent dynamics and slow purification}
\label{sec:tangent}

Exact occupation projectors obstruct a global invertible transfer coordinate, but they do not obstruct linear response about a physical trajectory.  In this section we construct the normalized tangent cocycle, extract its wall-localized slow sector, and connect that sector to purification from a maximally mixed state.

For a realized nonzero-probability branch, the normalized covariance update defines a map on its physical stratum,
\begin{equation}
 G_{t+1}=\Phi_{\xi_t}[G_t].
 \label{eq:normalized-map}
\end{equation}
Linearizing this map propagates a perturbation without inverting the Kraus operator,
\begin{equation}
 \delta G_{t+1}
 =D\Phi_{\xi_t}[G_t](\delta G_t)
 =\cJ_{t,\xi}^{\dagger}\,\delta G_t\,\cJ_{t,\xi},
 \label{eq:tangent-map}
\end{equation}
where the final factorization is taken on the support that survives the realized branch.  Composing these differentials gives the chronological one-leg cocycle
\begin{equation}
 \cT_{T,\xi}=\cJ_{0,\xi}\cJ_{1,\xi}\cdots\cJ_{T-1,\xi}.
 \label{eq:tangent-product}
\end{equation}
with the first-applied factor written leftmost.  The product may lose rank, but its finite-time singular values and forward images remain well defined.  We therefore use a stabilized forward filtration, which reduces to the familiar Oseledec/Lyapunov construction when the factors are invertible~\cite{Oseledec1968,Ludwig2013}, without assuming backward subspaces through directions that have been projected out.  The finite-time one-leg decay rates are
\begin{equation}
 \gamma_a(T)=-\frac{1}{T}\log s_a(\cT_{T,\xi}),
 \label{eq:lyapunov}
\end{equation}
using a sign convention in which small positive $\gamma_a$ denotes a slowly decaying one-leg direction.  Because the covariance tangent is a congruence, its physical decay rates are the allowed pair sums $\gamma_a+\gamma_b$, rather than the individual one-leg rates.

The maximally mixed reference $G_0=0$ makes the connection to purification especially direct.  For the full start-to-$T$ cocycle accumulated along that same maximally mixed record, the exact endpoint relation
\begin{equation}
 \cT_{T,\xi}^{\dagger}\cT_{T,\xi}=4C_T(\Id-C_T)
 \label{eq:maxmix-tangent-relation}
\end{equation}
matches the tangent singular values to the remaining mixed occupations.  This relation does not identify a separately sampled stationary pure-start QR spectrum.  A finite one-leg lower edge $\gamma_*>0$ produces exponential late-time purification.  If the matched wall cocycle instead has a low-rate density $\rho(\gamma)\sim\gamma^\alpha$, the entropy deficit scales as
\begin{equation}
 \Delta S_{\mathrm{pur}}(t)
 \sim\int_0^\infty\dd\gamma\,\rho(\gamma)e^{-2\gamma t}
 \sim t^{-(\alpha+1)}.
 \label{eq:purification-density}
\end{equation}
Figure~\ref{fig:tangent}(d) gives \ResultValue{$\alpha$ AND PARAMETER-FREE PURIFICATION EXPONENT}.  The exponent extracted from the tangent density is compared with Fig.~\ref{fig:state-cft}(d) without independently refitting the purification curve to force agreement.

We now characterize the slow tangent sector itself.  After taking the long-time limit at each fixed transverse size, Fig.~\ref{fig:tangent}(b) gives \ResultValue{WALL-ON GAP SCALING AND FINITE WALL-OFF GAP}.  Because several slow directions can become nearly degenerate, Fig.~\ref{fig:tangent}(c) displays the projector density of the complete low-exponent subspace rather than a single QR vector, yielding \ResultValue{WALL WEIGHT AND LOCALIZATION LENGTH}.  The trajectory-conditioned spectrum is finally compared with the unconditional covariance map in Fig.~\ref{fig:tangent}(e), where we find \ResultValue{CONDITIONAL--UNCONDITIONAL GAP COMPARISON}.  These tests place the present rank-losing circuit alongside monitored Lyapunov constructions while keeping the branch tangent distinct from an invertible logarithmic generator~\cite{Oshima2025,XiaoKawabata2026}.

For visualizing the finite-time singular directions, we also use the Hermitian log-stretch generator
\begin{equation}
 P_{T,\xi}^{\Lambda}=\frac{1}{2T}\log(\cT_{T,\xi}\cT_{T,\xi}^{\dagger})
\label{eq:quenched-log-object}
\end{equation}
on the surviving support.  The object $P_{T,\xi}^{\Lambda}$ is not a density matrix and need not be positive.  Its eigenvectors are the singular directions of one finite-time trajectory product, not eigenvectors obtained by averaging trajectory vectors.  Any quenched average is therefore applied either to this logarithmic object or to continuously tracked exponents, with the order of operations stated explicitly.

The same singular structure admits a scattering-network interpretation.  A regulated Gaussian Choi relation separates into an invertible quotient core and occupation caps that pin exactly filled or empty directions, while fresh gain/loss ancillas act as physical reservoir ports.  Chronological composition eliminates internal ports by a Schur complement, and the caps retain the rank information discarded by an ordinary transfer inverse.  Removing the regulator leaves the normalized tangent core in Eq.~\eqref{eq:tangent-map}; Appendix~\ref{app:cap-core-anchor} gives the cap bookkeeping, particle/hole symmetry, and regular-transfer limit.

\begin{figure*}[t]
  \placeholderfigure[0.97\textwidth]{
    FIGURE 5 PLACEHOLDER\\[0.5em]
    (a) Stabilized tangent-product workflow for exact projectors and gain/loss branches.\\[0.35em]
    (b) Wall-on and wall-off minimum tangent gaps versus transverse size, including zero-intercept and alternative finite-size fits.\\[0.35em]
    (c) Globally normalized projector density of the complete near-marginal subspace, averaged over Born trajectories.\\[0.35em]
    (d) Low-exponent density and parameter-free comparison with the purification curve of Fig.~\ref{fig:state-cft}.\\[0.35em]
    (e) Trajectory-conditioned result versus the mixed unconditional covariance/channel control.}
  \caption{
  Tangent mechanism for slow wall purification.  (a) Stabilized product of normalized branch differentials.  (b) Long-time-extrapolated wall-on and wall-off tangent gaps versus transverse size.  (c) Basis-independent projector density of the complete near-degenerate slow subspace, averaged over Born trajectories.  (d) Low-exponent density and its parameter-free prediction for purification.  (e) Comparison with the unconditional covariance map.  Component-resolved plots preserve their relative physical weights.  \ResultPlaceholder{Call the tangent sector thermodynamically gapless only after the $T\to\infty$ extrapolation at each fixed size and the subsequent zero-intercept conclusion are stable to fit form, aspect ratio, initial tangent frame, and trajectory resampling; otherwise use ``finite-size closing.''}  \ResultPlaceholder{Accept the mechanistic purification claim only if the same-record low-exponent density predicts the late-time decay without an independently adjustable exponent.}}
  \label{fig:tangent}
\end{figure*}

\section{Robustness and ensemble controls}
\label{sec:controls}

We finally test whether the wall signatures persist away from the ideal perfect-correction limit.  We introduce \ResultValue{IMPLEMENTED IMPERFECTION} and follow the bulk marker, wall entropy coefficient, tangent gap, and one orientation-sensitive response [Appendix Fig.~\ref{fig:app-robustness}].  This analysis gives \ResultValue{ROBUSTNESS INTERVAL OR THRESHOLD SCALING}.  Feedback-driven entanglement, steering, control, and absorbing-state transitions are already known in interacting and random adaptive circuits~\cite{IadecolaGaneshanPixley2023,RavindranathHanYang2023,SierantTurkeshi2023,ODeaMorningstarGopalakrishnan2024,PanGaneshanIadecola2024}; transport-limited frustration-free control provides a complementary preparation problem~\cite{Doerstel2026}.  Our purpose here is the stability of the same topological defect; a complete absorbing-transition analysis would additionally require an order parameter, finite-size crossings, and a dynamical exponent.  \ResultPlaceholder{Claim robustness only if this interval remains nonzero under size scaling; a single-size plateau is not evidence for a stable phase.}

The earlier Bhuiyan--Pan--Jian noise study suggests a natural endpoint for this robustness window~\cite{BhuiyanPanJian2026}.  Once noise destroys the bulk dynamical topological distinction or closes the relevant purity gap, a critical wall is no longer topologically required, although accidental boundary criticality is not excluded.  The measured relation between the two scales is \ResultValue{BULK--WALL THRESHOLD RELATION}.  A boundary crossover or preemptive wall transition can therefore occur before the bulk threshold; coincidence is not imposed in the analysis.  \ResultPlaceholder{State coincidence of wall and bulk thresholds only if independent finite-size crossings and uncertainty estimates resolve a common critical point; otherwise report separate crossover or transition scales.}

A matched postselected control addresses a different question.  For the same local schedule, Appendix~\ref{app:ensemble-distinctions} compares Born sampling with a forced favorable record and, where meaningful, an explicitly designed imaginary-time filter.  The initial overlap, Trotter step, normalization, and success probability are retained so that the comparison isolates which wall and record observables require the physical Born ensemble.

The averaged channel provides the final ensemble control.  Averaging Gaussian Kraus trajectories need not produce a Gaussian many-body channel, so we use closed covariance formulas only where closure is explicitly demonstrated.  On that well-defined sector, Fig.~\ref{fig:tangent}(e) gives \ResultValue{AVERAGED-CHANNEL GAP OR OCCUPATION-BRANCH RESULT}.  A half-occupied mode or linear occupation branch is reported as a property of the mixed channel and is not identified with the trajectory CFT.  \ResultPlaceholder{Any statement that criticality is ``hidden in trajectories'' must be supported by a matched channel calculation showing the absence of the same finite-size singularity in a well-defined channel observable.}

\section{Conclusion and outlook}
\label{sec:conclusion}

\subsection{Summary of results}
\label{sec:discussion}

In this work, we have used the perfect-correction limit of an adaptive Chern circuit to investigate the universal dynamics of a programmable domain wall.  Born sampling does not prepare one corrected wave function: it generates a charge-fluctuating ensemble of inequivalent Gaussian states.  Across that ensemble, the wall exhibits \ResultValue{FINAL STATE-CFT RESULT}, while the modular, response, and twist probes establish \ResultValue{FINAL CHIRALITY RESULT}.  The two microscopic wall constructions and the exact-Hamiltonian benchmarks separate these boundary properties from details of the finite-support controller.

The singular tangent formalism supplies the dynamical mechanism.  The differential of the normalized covariance map remains meaningful through rank-changing occupation projectors and gives \ResultValue{FINAL TANGENT AND PURIFICATION RESULT}.  The replicated Born weights independently yield \ResultValue{FINAL RECORD RESULT}, with the state and record central charges kept as distinct quantities.  Comparison with the unconditional covariance map gives \ResultValue{FINAL CHANNEL COMPARISON}.  Together, these observations define a three-layer bulk--boundary structure: a monitored class-A bulk, a chiral critical wall in its physical Born trajectories, and a defect theory in the probability distribution of those trajectories.  \ResultPlaceholder{Final synthesis paragraph to be rewritten around the subset of state, chirality, tangent, record, and channel claims that survive the complete acceptance table.}

This structure distinguishes the present problem from adaptive preparation of a common target state.  Local measurement and feedforward can compile conformal and chiral states, including through parton resources and multiscale adaptive circuits~\cite{LuLessaKimHsieh2022}.  Here those ingredients instead act repeatedly and locally, leaving the realized record and total charge stochastic.  The earlier Bhuiyan--Pan--Jian construction supplies the postselection-free bulk controller, programmable wall, and slow-purification observation~\cite{BhuiyanPanJian2026}; the additional result is the identification of the wall theory across state, orientation-sensitive, tangent, record, and channel observables.

The tangent and record constructions likewise extend, rather than replace, existing monitored-fermion descriptions.  Gaussian tensor-network and transfer approaches relate monitored dynamics to localization, and Lyapunov classifications identify topological boundary sectors~\cite{Jian2022,PanShapourianJian2025,XiaoKawabata2026,Oshima2025}.  Our normalized cocycle carries this analysis through exact rank loss and connects its low-exponent density quantitatively to purification.  Non-Hermitian or no-click Hamiltonians remain useful descriptions of selected branches, dissipative preparation constrains the unconditional channel, and continuous counting theory characterizes physical stochastic records~\cite{Gong2018,Goldstein2019,Landi2024,Starchl2025}.  The present contribution is their conjunction at a programmable two-dimensional Chern interface, together with a replica-calibrated defect construction rather than the bare relation $Z_\xi=p_\xi$.  Monitored free-fermion criticality supplies both the nearby phenomena and the necessary large-scale crossover cautions~\cite{AlbertonBuchholdDiehl2021,FavaPiroliSwann2023,ChahineBuchhold2024,NiedereggerVovkStarchl2026}.  Interactive feedback and control provide a distinct setting in which entanglement and absorbing transitions can coexist~\cite{IadecolaGaneshanPixley2023,ODeaMorningstarGopalakrishnan2024,PanGaneshanIadecola2024}.

\subsection{Outlook and future directions}
\label{sec:outlook}

The first extension is to interacting adaptive circuits.  Covariance closure is then lost, but the distinction among a selected branch, the Born ensemble, the record distribution, and the unconditional channel remains.  Determining which parts of the tangent and record constructions survive beyond Gaussian dynamics would provide a route toward interacting chiral defect theories.

A second direction is the competition between adaptive correction and noise or an absorbing process.  The robustness analysis here follows one imperfection across the bulk marker and wall observables.  A complete transition theory would determine whether the wall loses criticality at the bulk topological threshold, crosses over earlier, or meets a separate absorbing or control transition, with independent order parameters and dynamical scaling.

Finally, the fresh-ancilla realization gives a direct quantum-simulation target.  Refreshed occupied and empty modes, midcircuit occupation measurements, local feedforward, and fermionic exchange operations implement the ideal update using the BPJ route.  The experiment would prepare and sample a universal charge-fluctuating ensemble with a programmable defect, rather than certify one Chern wave function.  Its most distinctive observable may be the comparison between universal quantities extracted from the classical record and those requiring trajectory-state tomography.  Establishing the relative sample and hardware cost of those two routes is an important operational benchmark for future native-fermion processors~\cite{GonzalezCuadraBluvsteinKalinowski2023,CalliariFromonteilCesa2026,OttGonzalezCuadraZache2025,SchuckertCraneGorshkov2024}.

\begin{acknowledgments}
We thank collaborators and colleagues to be listed in the final manuscript.  Funding acknowledgments will be added before submission.
\end{acknowledgments}

\appendix

\section{Microscopic instrument and exact perfect correction}
\label{app:microscopic-instrument}

This appendix fixes the microscopic convention used throughout the paper.  In
the main text, a controller channel carries the BPJ label
\(j=(\bfr,\nu,\sigma)\), with \(\nu=A,B\), \(\sigma=\pm\), localized operator
\(\h\chi_{\bfr,\nu,\sigma}\), and occupation
\(\hcN_{\bfr,\nu,\sigma}=\h\chi_{\bfr,\nu,\sigma}^\dagger
\h\chi_{\bfr,\nu,\sigma}\).  The coefficient column of that operator is
denoted by the corresponding plain symbol \(\chi_j\) below.  Thus the generic
notation \(a_\chi\) and \(n_\chi\) used in the derivations is precisely
\(\h\chi_j\) and \(\hcN_j\) for a selected BPJ channel, not an additional
fermion mode.  We consider \(N\) complex fermion modes with
\begin{equation}
    \{c_i,c_j^\dagger\}=\delta_{ij},
    \qquad
    \{c_i,c_j\}=0,
\end{equation}
and use the number-conserving two-point matrix and centered covariance
\begin{equation}
    C_{ij}=\langle c_i^\dagger c_j\rangle,
    \qquad
    G=2C-\Id .
    \label{eq:app-C-G-convention}
\end{equation}
Thus \(0\le C\le\Id\), while the eigenvalues of \(G\) lie in
\([-1,1]\).  A pure Slater determinant obeys \(C^2=C\), equivalently
\(G^2=\Id\).  All formulas below act on the physical one-particle space; an
orbital projector is denoted by \(P_\chi=\chi\chi^\dagger\), with
\(\chi^\dagger\chi=1\).

\subsection{Occupation measurement and branch probabilities}

For the normalized orbital
\begin{equation}
    a_\chi=\sum_j\chi_j c_j,
    \qquad
    n_\chi=a_\chi^\dagger a_\chi,
\end{equation}
the two projectors are
\begin{equation}
    \Pi_1=n_\chi,
    \qquad
    \Pi_0=\Id-n_\chi=a_\chi a_\chi^\dagger.
    \label{eq:app-occupation-projectors}
\end{equation}
It is convenient to write \(\eta_m=2m-1\), so that \(\eta_m=+1\)
for an occupied result and \(\eta_m=-1\) for an empty result.  On a
Gaussian state the Born probabilities are
\begin{equation}
    p_m(G)=\Tr(\rho\Pi_m)
    =\frac{1+\eta_m g_\chi}{2},
    \qquad
    g_\chi=\chi^\dagger G\chi.
    \label{eq:app-binary-born-probability}
\end{equation}
Only outcomes with \(p_m>0\) define normalized branches.  A formal pole at
\(p_m=0\) is therefore a zero-probability branch, not a failure of the
Gaussian update.

Let \(s\in\{0,1\}\) be the target occupation assigned by the controller.
For the upper-band controller \(s=0\), while for the lower-band controller
\(s=1\).  Perfect correction is represented by the two system Kraus
operators
\begin{equation}
    \begin{array}{c|cc}
      &m=0&m=1\\ \hline
      s=0&K_{0|0}=\Pi_0&K_{1|0}=a_\chi\\[1mm]
      s=1&K_{0|1}=a_\chi^\dagger&K_{1|1}=\Pi_1
    \end{array}.
    \label{eq:app-perfect-correction-kraus-table}
\end{equation}
These operators satisfy, branch by branch,
\begin{align}
    K_{m|s}^\dagger K_{m|s}&=\Pi_m,
    &\Pi_sK_{m|s}&=K_{m|s},\nonumber\\
    \sum_mK_{m|s}^\dagger K_{m|s}&=\Id.
    \label{eq:app-perfect-correction-checks}
\end{align}
The first identity shows that correction does not alter the Born probability.
The second shows that the measured orbital has the target occupation after
every nonzero branch.  Randomness remains in the conditional state of the
overlapping orbitals; perfect correction does not collapse the full circuit
to a deterministic wavefunction.

\subsection{Fresh-ancilla dilation and deterministic fSWAP correction}

Introduce one fresh ancillary mode \(d\), and define
\begin{align}
    U_{\mathrm{fSWAP}}
    &=\fSWAP(a_\chi,d)
      =\exp\!\left[
      \ii\frac{\pi}{2}(a_\chi^\dagger-d^\dagger)(a_\chi-d)\right]
      \nonumber\\
    &=\Id-n_\chi-n_d+a_\chi^\dagger d+d^\dagger a_\chi.
    \label{eq:app-fswap-unitary}
\end{align}
It leaves \(\lvert00\rangle\) invariant, exchanges
\(\lvert01\rangle\) and \(\lvert10\rangle\), and gives the fermionic
exchange sign on \(\lvert11\rangle\).  After the occupation measurement,
the physical mode is already in the definite state \(\lvert m\rangle_\chi\).
If \(m\ne s\), the ancilla is prepared in \(\lvert s\rangle_d\) and
fSWAP is applied.  Consequently,
\begin{equation}
    \lvert m\rangle_\chi\lvert s\rangle_d
    \xrightarrow{\ U_{\mathrm{fSWAP}}\ }
    \lvert s\rangle_\chi\lvert m\rangle_d,
    \qquad m\ne s.
    \label{eq:app-fswap-deterministic-action}
\end{equation}
The outgoing ancilla has a known occupation and factorizes from the physical
system.  It may therefore be discarded, or locally reset and reused, without
postselection.  Taking the appropriate ancilla matrix element reproduces the
odd entries in Eq.~\eqref{eq:app-perfect-correction-kraus-table},
\begin{equation}
    {}_d\!\langle0|U_{\mathrm{fSWAP}}|1\rangle_d=a_\chi^\dagger,
    \qquad
    {}_d\!\langle1|U_{\mathrm{fSWAP}}|0\rangle_d=a_\chi,
    \label{eq:app-fswap-matrix-elements}
\end{equation}
but these matrix elements are a derivation of the reduced Kraus map, not an
instruction to postselect the ancilla.  Total charge is conserved on the
enlarged system--ancilla Hilbert space.

\subsection{Mixed fresh ancilla and Markovian bath control}
\label{app:mixed-fresh-ancilla}

A distinct control replaces the target Fock-state ancilla by a freshly drawn
mixed mode
\begin{equation}
    \rho_a(\bar n_a)
    =(1-\bar n_a)\lvert0\rangle\!\langle0\rvert_a
      +\bar n_a\lvert1\rangle\!\langle1\rvert_a,
    \qquad 0\leq \bar n_a\leq1 .
    \label{eq:app-mixed-ancilla-state}
\end{equation}
Apply the same fSWAP as in Eq.~\eqref{eq:app-fswap-unitary}, now between the
system orbital selected by \(P=P_\chi\) and this ancilla, and trace out the
outgoing ancilla.  With \(Q=\Id-P\), the reduced system covariance is
\begin{equation}
    \boxed{
    \mathcal B_{\bar n_a}(G)
    =QGQ+(2\bar n_a-1)P .}
    \label{eq:app-mixed-ancilla-channel}
\end{equation}
Thus the selected system mode is replaced by the bath occupation and all of
its incoming correlations with the complementary modes leave with the
discarded ancilla.

Equivalently, one may first sample the classical preparation label
\(a\in\{0,1\}\) of Eq.~\eqref{eq:app-mixed-ancilla-state}.  Writing
\(\eta_a=2a-1\), the two conditional replacement maps are
\begin{align}
    G'_+&=QGQ+P,
    &p_+&=\bar n_a,\nonumber\\
    G'_-&=QGQ-P,
    &p_-&=1-\bar n_a,
    \label{eq:app-mixed-ancilla-unravelling}
\end{align}
whose classical average is Eq.~\eqref{eq:app-mixed-ancilla-channel}.  These
branches are generally mixed even for a pure incoming system, because tracing
the outgoing ancilla discards any initial entanglement of the selected mode.
They remain pure only when the input is pure and the selected mode is already
uncorrelated with its complement, \(QGP=0\).  Alternatively, conditioning on
an occupation measurement of the outgoing ancilla restores a pure-state
measurement trajectory but gives the nonlinear conditional map of
Eq.~\eqref{eq:app-perfect-correction-G-map}, not
Eq.~\eqref{eq:app-mixed-ancilla-unravelling}.

Equation~\eqref{eq:app-mixed-ancilla-channel} is therefore a Markovian bath or
noise control, not the production perfect-correction rule: the latter chooses
the definite target occupation \(s\) after observing the system outcome.
Fresh preparation (or a complete reset) is essential.  Reusing the outgoing
ancilla without resetting it feeds the previous system mode back into later
steps and produces a non-Markovian instrument.

\appendixplaceholderfigure{fig:app-instrument}{Perfect-correction instrument and its local dilation}{
Panel (a): occupation measurement of a finite-support orbital, with the two
Born outcomes.  Panel (b): the four entries of
Eq.~\eqref{eq:app-perfect-correction-kraus-table}.  Panel (c): deterministic
gain and loss realized by an occupied or empty fresh ancilla followed by
fSWAP.  The final artwork should distinguish the retained physical record
from a postselected ancilla port and should indicate that the outgoing
ancilla is factorized and recyclable.}

\subsection{Conditional covariance map}

Let \(P=P_\chi\), \(Q=\Id-P\), and retain only a nonzero outcome \(m\).
The conditional measurement followed by perfect reset to \(s\) acts on the
centered covariance as
\begin{equation}
    \boxed{
    \Phi_{m|s}(G)
    =Q\bigl(\Id+\eta_m GP\bigr)^{-1}GQ+\eta_sP .}
    \label{eq:app-perfect-correction-G-map}
\end{equation}
The inverse is only a rank-one resolvent.  The Sherman--Morrison form makes
the probability denominator explicit,
\begin{equation}
    \bigl(\Id+\eta_m GP\bigr)^{-1}
    =\Id-\frac{\eta_m GP}{1+\eta_m g_\chi}.
    \label{eq:app-rank-one-resolvent}
\end{equation}
Equation~\eqref{eq:app-perfect-correction-G-map} applies equally to pure and
mixed Gaussian inputs.  The last term is set by the target, whereas the
state-dependent contraction of the orthogonal complement is selected by the
actual measurement result.

For reference, the unconditional perfect-correction channel for this one
orbital is especially simple.  Summing the unnormalized branches gives the
replacement channel
\begin{equation}
    \mathcal E_s(\rho)
    =\sum_mK_{m|s}\rho K_{m|s}^\dagger
    =\lvert s\rangle\!\langle s\rvert_\chi
      \mathbin{\widehat\otimes}\Tr_\chi\rho,
    \label{eq:app-reset-channel-density}
\end{equation}
where \(\widehat\otimes\) denotes the graded fermionic product.  Its affine
covariance action is
\begin{equation}
    \mathbb E_m\!\left[\Phi_{m|s}(G)\right]
    =QGQ+\eta_sP.
    \label{eq:app-reset-channel-G}
\end{equation}
This form is used on the parity-even, number-conserving states of the present
protocol; the graded tensor factorization is part of the statement.
The nonlinear conditioned map and the linear averaged channel must not be
interchanged in trajectory observables.

For an ordered record
\(\xi=(m_1,\ldots,m_D)\), with prescribed targets \(s_t\) and orbitals
\(\chi_t\), define
\begin{equation}
    G_t=\Phi_{m_t|s_t}(G_{t-1}).
\end{equation}
The complete Born weight is the chronological product
\begin{equation}
    \boxed{
    p_\xi(G_0)
    =\prod_{t=1}^{D}
      \frac{1+\eta_{m_t}\chi_t^\dagger G_{t-1}\chi_t}{2}.}
    \label{eq:app-record-born-weight}
\end{equation}
All later factors are evaluated on the state produced by the earlier
outcomes.  Evaluating every probability at \(G_0\), or attaching additional
weights to already Born-sampled trajectories, would give a different
ensemble.

The chronological channel convention is also part of the model.  At each
active unit cell, the manuscript follows BPJ and visits
\begin{align}
    \mathcal Q_{\rm BPJ}
    &=\bigl((A,-),(A,+),(B,-),(B,+)\bigr),
    \nonumber\\
    (s_q)_{q\in\mathcal Q_{\rm BPJ}}&=(1,0,1,0).
    \label{eq:app-bpj-channel-order}
\end{align}
Equivalently, the circuit fills the lower OW mode and then depletes the upper
OW mode for \(\nu=A\), before repeating the same pair for \(\nu=B\).
Because overlapping occupations do not generally commute, a word with the
same four labels in another order defines a different finite-range
instrument and cannot be identified with Eq.~\eqref{eq:app-bpj-channel-order}
by relabeling its output.

\section{Finite-support controllers and programmed domain walls}
\label{app:finite-support-controller}

This appendix supplies the Fourier conventions, the effect of real-space
windowing, and the implementation details suppressed in Sec.~\ref{sec:model}.
The distinction among the target-band Chern number, the winding of a
windowed mode's overlap with that band, and the Chern number of the compact
mode itself is particularly important.

\subsection{Uniform parent and projected frame}

On a finite periodic lattice, the real-space band projector associated with
Eq.~\eqref{eq:parent-bloch-projectors} is
\begin{equation}
    P_b(\mathbf r-\mathbf r';\alpha)
    =\frac{1}{N_xN_y}\sum_{\mathbf k}
      e^{i\mathbf k\cdot(\mathbf r-\mathbf r')}
      P_b(\mathbf k;\alpha).
    \label{eq:app-real-space-band-projector}
\end{equation}
For the orthonormal cell-local trials
\(\tau_A=(1,1)^{\mathsf T}/\sqrt2\) and
\(\tau_B=(1,-1)^{\mathsf T}/\sqrt2\), define the conventional projected
one-particle OW vector.  The appendix uses a generic band label \(b=\pm\);
it is the sign \(\sigma=\pm\) in the BPJ channel label of the main text:
\begin{equation}
    \lvert\widetilde W_{\mathbf R,\nu,b}\rangle
    =P_b(\alpha)\lvert\mathbf R,\tau_\nu\rangle.
    \label{eq:app-projected-frame-vector}
\end{equation}
The corresponding unnormalized annihilation operator is
\begin{align}
    \hat{\widetilde\chi}_{\mathbf R,\nu,b}
    &=\sum_{\mathbf r,\mu}
      \langle\widetilde W_{\mathbf R,\nu,b}\vert
      \mathbf r,\mu\rangle c_{\mathbf r\mu}
      \nonumber\\
    &=\frac{1}{\sqrt{N_xN_y}}\sum_{\mathbf k}
      e^{i\mathbf k\cdot\mathbf R}
      \tau_\nu^\dagger P_b(\mathbf k;\alpha)\mathbf c_{\mathbf k}.
    \label{eq:app-ow-operator-fourier}
\end{align}
The normalized BPJ operator is
\begin{equation}
    \h\chi^{(\infty)}_{\mathbf R,\nu,b}
    =\frac{\hat{\widetilde\chi}_{\mathbf R,\nu,b}}
    {\|\widetilde W_{\mathbf R,\nu,b}\|}.
    \label{eq:app-normalized-ow-operator}
\end{equation}
Writing \(N_c=N_xN_y\), the vector stored by the numerical covariance code is
the coefficient column
\begin{equation}
    \widetilde\chi_{\mathbf R\nu b}(\mathbf r,\mu)
    =\frac{1}{N_c}\sum_{\mathbf k}
      e^{-i\mathbf k\cdot(\mathbf r-\mathbf R)}
      \bigl[P_b(\mathbf k;\alpha)\tau_\nu\bigr]^*_{\mu}.
    \label{eq:app-stored-coefficient-convention}
\end{equation}
It equals
\(\langle\widetilde W_{\mathbf R,\nu,b}\vert\mathbf r,\mu\rangle\).
This conjugation is why the filled-lower-band covariance in the convention
\(C_{ij}=\langle c_i^\dagger c_j\rangle\) is
\(C_{\rm CI}(\mathbf k)=P_-^*(\mathbf k)\).  The physical one-particle
projector is recovered as \(C_{\rm CI}^{\mathsf T}=P_-\); hence the target
Chern number remains \(+1\) when the same convention is used in the marker.
The transpose must be retained when comparing analytic one-particle kets with
code covariance projectors.

The word ``overcomplete'' is made precise by the unnormalized tight-frame
identity
\begin{align}
    \sum_{\mathbf R,\nu}
      \lvert\widetilde W_{\mathbf R,\nu,b}\rangle
      \langle\widetilde W_{\mathbf R,\nu,b}\rvert
    &=P_b\left(\sum_{\mathbf R,\nu}
      \lvert\mathbf R,\tau_\nu\rangle
      \langle\mathbf R,\tau_\nu\rvert\right)P_b
      \nonumber\\
    &=P_b.
    \label{eq:app-tight-frame-proof}
\end{align}
There are two projected vectors per cell for a one-band subspace of one state
per cell.  They span that subspace without constituting an orthogonal Wannier
basis.  After normalizing every nonzero full-support mode, the filled lower
band obeys the common target relations in
Eq.~\eqref{eq:ow-stabilizer-target}.

\subsection{Atomic large-mass limit}
\label{app:atomic-large-mass-limit}

The large positive trivial mass has a simple gauge-independent limit.  For
the Bloch vector in Eq.~\eqref{eq:parent-bloch-projectors},
\begin{equation}
    \mathbf d(\mathbf k;\alpha)
    =\bigl(\sin k_x,\sin k_y,
      \alpha-\cos k_x-\cos k_y\bigr),
\end{equation}
the spectral projectors may be written
\begin{equation}
    P_\pm(\mathbf k;\alpha)
    =\frac12\left[\Id\pm
      \widehat{\mathbf d}(\mathbf k;\alpha)\cdot\boldsymbol\sigma\right].
    \label{eq:app-projector-unit-vector}
\end{equation}
At fixed momentum,
\begin{equation}
    \widehat{\mathbf d}(\mathbf k;\alpha)
    =\widehat{\mathbf z}+O(\alpha^{-1}),
    \qquad
    P_\pm(\mathbf k;\alpha)
    =\frac12(\Id\pm\sigma_z)+O(\alpha^{-1})
    \label{eq:app-large-mass-projectors}
\end{equation}
as \(\alpha\to+\infty\).  The leading correction is momentum dependent and
off diagonal in the orbital basis, while the diagonal correction begins at
higher order.

Fourier transforming Eq.~\eqref{eq:app-large-mass-projectors} gives
\begin{equation}
    P_\pm(\mathbf r-\mathbf r';\alpha)
    =\delta_{\mathbf r,\mathbf r'}
      \frac{\Id\pm\sigma_z}{2}+O(\alpha^{-1}).
    \label{eq:app-atomic-projector-kernel}
\end{equation}
Consequently the normalized projected trial vectors of
Eq.~\eqref{eq:app-projected-frame-vector} collapse onto the corresponding
onsite orbital whenever the trial has nonzero overlap with that orbital.  The
choice \(\alpha_{\rm triv}=30\) is therefore a controlled proxy for an atomic
trivial exterior.  The support-terminated controller of
Eq.~\eqref{eq:app-support-terminated-interface} is nevertheless a separate
active boundary condition: masking a topological OW frame and omitting
exterior centers is not identical to evolution with any finite value of
\(\alpha_{\rm triv}\), even though the infinite-mass limit motivates the
hard-boundary interpretation.

\subsection{Truncate--then--renormalize rule}

For integer shell radius \(n\), let
\begin{equation}
    M_{\mathbf R}^{(n)}(\mathbf r)
    =\mathbf 1\!\left[
      |d_{N_x}(x,R_x)|\leq n\ \text{and}\
      |d_{N_y}(y,R_y)|\leq n\right].
    \label{eq:app-square-support}
\end{equation}
The measured operator coefficient is obtained by applying every geometric
mask and only then normalizing,
\begin{equation}
    \chi^{(n)}_{\mathbf R\nu b}
    =\frac{M_{\mathbf R}^{(n)}\widetilde\chi_{\mathbf R\nu b}}
    {\|M_{\mathbf R}^{(n)}\widetilde\chi_{\mathbf R\nu b}\|},
    \qquad
    a^{(n)}_{\mathbf R\nu b}
    =\sum_{\mathbf r,\mu}
      \chi^{(n)}_{\mathbf R\nu b}(\mathbf r,\mu)c_{\mathbf r\mu}.
    \label{eq:app-finite-support-frame}
\end{equation}
Thus \(n=1\) is a \(3\times3\)-cell square with at most 18 microscopic
modes.  The special half-shell option is instead the five-cell cross
\(\{(0,0),(\pm1,0),(0,\pm1)\}\).  Normalization is performed independently
for every \((\mathbf R,\nu,b)\).  Windowing introduces out-of-band weight,
so the compact occupations cease to be a mutually compatible set of exact
stabilizers even though the controller remains strictly local.

\subsection{Windowed overlap and the surviving obstruction}

Let \(\lvert u_b(\mathbf k)\rangle\) be a normalized Bloch eigenstate and
\(f_{\nu b}(\mathbf k)=\langle u_b(\mathbf k)\vert\tau_\nu\rangle\) its
trial-state form factor.  Multiplication by a real-space window becomes a
twisted convolution.  For a mode centered at the origin, its overlap with
the target band is
\begin{equation}
    f^{(n)}_{\nu b}(\mathbf k')
    =\frac{1}{N_xN_y}\sum_{\mathbf k}
      \widetilde M_n(\mathbf k'-\mathbf k)
      f_{\nu b}(\mathbf k)
      \langle u_b(\mathbf k')\vert u_b(\mathbf k)\rangle,
    \label{eq:app-windowed-twisted-convolution}
\end{equation}
where the square window has
\begin{align}
    \widetilde M_n(\mathbf q)&=D_n(q_x)D_n(q_y),
    \nonumber\\
    D_n(q)&=\sum_{j=-n}^{n}e^{-iqj}
    =\frac{\sin[(n+\tfrac12)q]}{\sin(q/2)}.
    \label{eq:app-square-dirichlet-kernel}
\end{align}
Restoring a nonzero center adds only the translation phases displayed in the
full derivation of the windowed controller.  Dividing the compact mode by its
real-space norm multiplies Eq.~\eqref{eq:app-windowed-twisted-convolution}
by a positive constant and therefore does not change its zeros or phase
winding.

For a small positively oriented loop \(\Gamma\) enclosing an isolated zero,
define
\begin{equation}
    \nu_{\rm ov}^{(n)}(\Gamma)
    =\frac{1}{2\pi i}\oint_\Gamma
      \dd\mathbf k\cdot\boldsymbol\nabla_{\mathbf k}
      \log f^{(n)}_{\nu,-}(\mathbf k).
    \label{eq:app-overlap-winding}
\end{equation}
For the unwindowed form factor, the convention of
Eq.~\eqref{eq:parent-bloch-projectors} gives
\(\mathcal C[P_-]=-\nu_{\rm ov}^{(\infty)}\).  At
\(\alpha=1\), the $A$ trial has its unwindowed zero at
\(\mathbf k_0/\pi=(0.5,0)\).  Direct evaluation of the square-window
convolution gives
\begin{equation}
    \begin{array}{c|c|c|c}
      n_{\rm shell}&\text{support}&
      \|M^{(n)}\widetilde W\|^2/\|\widetilde W\|^2&
      \mathbf k_*^{(n)}/\pi\\ \hline
      1&3\times3&0.992155&(0.4922,0)\\
      2&5\times5&0.999137&(0.4925,0)
    \end{array}
    \label{eq:app-windowed-overlap-results}
\end{equation}
For both rows, the descendant zero has \(\nu_{\rm ov}^{(n)}=-1\).  The zero
is displaced rather than filled.  Its local form remains
\(f^{(n)}(\mathbf k_*+\mathbf p)=A_n(p_x-ip_y)+O(p^2)\), which fixes the
reported winding.  Equation~\eqref{eq:app-windowed-overlap-results}, rather
than the retained norm alone, is the finite-window check relevant to the
controller.

For a reproducible discrete audit, sample a loop
\(\Gamma=(\mathbf k_0,\ldots,\mathbf k_{N_\Gamma}=\mathbf k_0)\) and compute
\begin{equation}
    \nu_{\rm loop}^{(n)}
    =\frac{1}{2\pi}\sum_{j=0}^{N_\Gamma-1}
      \operatorname{Arg}\!\left[
      \frac{f^{(n)}(\mathbf k_{j+1})}{f^{(n)}(\mathbf k_j)}\right].
    \label{eq:app-discrete-overlap-winding}
\end{equation}
The loop must avoid the zero and remain within a region containing no other
zeros.  The result is checked under increasing momentum resolution, several
loop radii, and smooth gauge changes of \(u_-(\mathbf k)\); the phase ratios
make the final integer gauge invariant when the loop is closed consistently.

The overlap winding is not the Chern number of the compact-support spinor.
Writing its Bloch transform as the two-component trigonometric polynomial
\begin{align}
    \lvert\psi_n(\mathbf k)\rangle
    &=\sum_{\mathbf r\in\operatorname{supp}M^{(n)}}
      e^{-i\mathbf k\cdot\mathbf r}C_{\mathbf r}\lvert\tau_\nu\rangle,
    \nonumber\\
    C_{\mathbf r}&=\int_{\rm BZ}\frac{\dd^2q}{(2\pi)^2}
      e^{i\mathbf q\cdot\mathbf r}P_-(\mathbf q),
    \label{eq:app-compact-spinor}
\end{align}
one sees that, whenever \(\psi_n(\mathbf k)\neq0\) throughout the Brillouin
torus, its normalized version is a smooth periodic global section.  Direct
minimization for the \(n=1\) mode gives
\(\min_{\mathbf k}\|\psi_1(\mathbf k)\|=5.106\times10^{-2}\), so this
condition is satisfied.  Hence
\begin{equation}
    \mathcal C_{\rm target}=+1,
    \qquad \nu_{\rm ov}^{(1)}=-1,
    \qquad \mathcal C_{\rm compact}^{(1)}=0.
    \label{eq:app-three-topological-objects}
\end{equation}
A coarse Fukui--Hatsugai--Suzuki discretization of an almost-vanishing compact
spinor can cross its admissibility bound and return a spurious integer.  We
therefore apply FHS to the reconstructed ket-space target projector
\(C_{\rm CI}^{\mathsf T}=P_-\), and evaluate real-space markers with the same
transpose/orientation convention.  The loop winding instead diagnoses the
overlap of that target with the compact controller; the two estimators are not
identified~\cite{FukuiHatsugaiSuzuki2005}.

\appendixplaceholderfigure{fig:app-windowed-controller}{Finite window and overlap obstruction}{
Panel (a): real-space weight of the unwindowed $A-$ mode and the $3\times3$
window.  Panel (b): $|f^{(n)}_{A,-}(k_x,0)|$ for the unwindowed, $n=1$, and
$n=2$ modes, marking the displaced zeros.  Panel (c): phase circulation on
loops surrounding the $n=1$ zero.  Panel (d): target-band FHS invariant,
overlap winding, and compact-spinor FHS result shown as three explicitly
different diagnostics.}

\subsection{Exact implementation of the two walls}

Let
\(\mathcal R_{\rm top}=\{\mathbf R:x_L\leq R_x\leq x_R\}\), with the
endpoints stated explicitly for every geometry.  The explicit-interface
circuit uses
\begin{align}
    \alpha_{\mathbf R}&=
    \begin{cases}
      \alpha_{\rm top},&\mathbf R\in\mathcal R_{\rm top},\\
      \alpha_{\rm triv},&\mathbf R\notin\mathcal R_{\rm top},
    \end{cases}
    \nonumber\\
    \chi^{(n),{\rm exp}}_{\mathbf R\nu b}
    &=\frac{M_{\mathbf R}^{(n)}\widetilde\chi_{\mathbf R\nu b}
      (\alpha_{\mathbf R})}
    {\|M_{\mathbf R}^{(n)}\widetilde\chi_{\mathbf R\nu b}
      (\alpha_{\mathbf R})\|}.
    \label{eq:app-explicit-interface}
\end{align}
All centers remain active.  Each center uses a projector from one of two
\emph{uniform} parents, and its compact tail may cross the interface.  It is
therefore inaccurate to describe Eq.~\eqref{eq:app-explicit-interface} as
the spectral frame of a single spatially inhomogeneous Hamiltonian.

The code constructs the more general same-region mask
\begin{equation}
    D_{\mathbf R}(\mathbf r)
    =\Theta_{\rm top}(\mathbf R)\Theta_{\rm top}(\mathbf r)
    +[1-\Theta_{\rm top}(\mathbf R)]
     [1-\Theta_{\rm top}(\mathbf r)].
    \label{eq:app-same-region-mask}
\end{equation}
The support-terminated protocol studied in the paper then restricts the
active set to \(\mathbf R\in\mathcal R_{\rm top}\), for which
\(D_{\mathbf R}(\mathbf r)=\Theta_{\rm top}(\mathbf r)\), and uses
\begin{equation}
    \chi^{(n),{\rm term}}_{\mathbf R\nu b}
    =\frac{D_{\mathbf R}M_{\mathbf R}^{(n)}
      \widetilde\chi_{\mathbf R\nu b}(\alpha_{\rm top})}
    {\|D_{\mathbf R}M_{\mathbf R}^{(n)}
      \widetilde\chi_{\mathbf R\nu b}(\alpha_{\rm top})\|},
    \qquad \mathbf R\in\mathcal R_{\rm top}.
    \label{eq:app-support-terminated-interface}
\end{equation}
Both \(D_{\mathbf R}\) and \(M_{\mathbf R}^{(n)}\) are diagonal real-space
multiplication operators, so they commute.  The order in
Eq.~\eqref{eq:app-support-terminated-interface} is merely notational: the
product mask defines the orbital before the circuit begins, and the result is
normalized once.  Exterior canonical orbitals are measured once to prepare a
Born-sampled onsite product state and are omitted from subsequent OW cycles.
The exterior outcome word \(\zeta\), its probability \(p_\zeta\), and the
conditional state \(\rho_{0|\zeta}\) are stored separately; the complete
trajectory weight is
\(p_{\zeta,\xi|\pi}=p_\zeta p_{\xi|\pi,\zeta}\).
A two-sided hard-truncation protocol that continues to update exterior
centers is a different circuit and is not pooled with the support-terminated
data.

On a periodic lattice the inclusive cell interval \([x_L,x_R]\) places the
physical interfaces on the bonds \(x_L-\tfrac12\) and
\(x_R+\tfrac12\).  Quoted wall widths must use these explicit endpoints;
implementation-dependent default fractions of \(N_x\) are not part of the
model definition.  Finally, the exact-$H$ benchmark is obtained by
diagonalizing a genuinely inhomogeneous single-particle Hamiltonian.  It
tests the expected finite-geometry boundary theory but is neither
Eq.~\eqref{eq:app-explicit-interface} nor
Eq.~\eqref{eq:app-support-terminated-interface}.


\section{Singular normalized tangent cocycle}
\label{app:tangent-cocycle}

The appropriate linear response of a conditioned trajectory is the
derivative of the normalized covariance map.  Differentiating
Eq.~\eqref{eq:app-perfect-correction-G-map} at a nonzero branch gives
\begin{equation}
    D\Phi_{m|s}\big|_G[H]
    =J_m(G)^\dagger HJ_m(G),
    \label{eq:app-tangent-congruence}
\end{equation}
with the half-Jacobian
\begin{equation}
    \boxed{
    J_m(G)
    =Q-\eta_m\frac{PGQ}{1+\eta_m g_\chi}.}
    \label{eq:app-projective-half-jacobian}
\end{equation}
The target \(s\) does not enter this derivative: correction replaces the
measured mode by a constant covariance, while the realized outcome controls
the response of its complement.  In particular,
\begin{equation}
    J_m(G)\chi=0.
    \label{eq:app-half-jacobian-null}
\end{equation}
Exact projectors therefore produce a semi-invertible cocycle.  No inverse of
the accumulated product is required.

For a normalized many-body map
\(\rho'=K\rho K^\dagger/p_K\), the same statement follows from the quotient
rule
\begin{equation}
    \delta\rho'
    =\frac{1}{p_K}\left[
      K\delta\rho K^\dagger
      -\rho'\Tr(K\delta\rho K^\dagger)
    \right].
    \label{eq:app-many-body-quotient-rule}
\end{equation}
The subtraction removes the norm-changing direction.  It is precisely this
Born normalization that distinguishes \(J_m(G)\) from a bare one-particle
transfer matrix.

\subsection{Chronological order and QR accumulation}

Along a fixed record, let \(J_t=J_{m_t}(G_{t-1})\).  Repeated use of
Eq.~\eqref{eq:app-tangent-congruence} yields
\begin{equation}
    \delta G_D
    =\widetilde J_D^\dagger\delta G_0\widetilde J_D,
    \qquad
    \boxed{\widetilde J_D=J_1J_2\cdots J_D.}
    \label{eq:app-chronological-half-cocycle}
\end{equation}
The first-applied factor is leftmost in \(\widetilde J_D\).  A column QR
implementation propagates the adjoint factors in the reverse-looking, but
equivalent, order:
\begin{equation}
    Y_t=J_t^\dagger Q_{t-1}=Q_tR_t,
    \qquad
    \widetilde J_t^\dagger Q_0
    =Q_tR_tR_{t-1}\cdots R_1.
    \label{eq:app-qr-order}
\end{equation}
This identity fixes the convention and prevents an accidental reversal of a
noncommuting word.

If QR is performed once per circuit cycle, the finite-time one-leg growth
rates in units of cycles are
\begin{equation}
    \ell_a^{(T)}(\xi)
    =\frac{1}{T}\sum_{t=1}^{T}\log |(R_t)_{aa}|.
    \label{eq:app-finite-time-lyapunov}
\end{equation}
The diagonal phases of \(R_t\) are absorbed into \(Q_t\), so the reported
diagonal is nonnegative.  A rank loss gives an exact zero and hence
\(-\infty\); a numerical implementation should report null or censored
counts separately rather than silently replacing them by a physical finite
rate.

The Oseledets interpretation requires more than the algebra above.  For each
finite lattice one first takes \(T\to\infty\) along a stationary ergodic
Born process, subject to the logarithmic integrability assumptions of the
multiplicative ergodic theorem~\cite{Oseledec1968}.  Only after this temporal
limit is controlled should one take the thermodynamic limit.  For \(S\)
independent trajectories the quenched estimator is
\begin{equation}
    \overline{\ell_a^{(T)}}
    =\frac{1}{S}\sum_{\xi=1}^{S}
      \frac{1}{T}\sum_t\log |(R_{t,\xi})_{aa}|.
    \label{eq:app-quenched-lyapunov-estimator}
\end{equation}
It is neither the logarithm of an averaged product nor the rate obtained from
the singular values of \(\mathbb E_\xi\widetilde J_T\).  When the leading
physical rates are nonpositive, the corresponding one-leg decay rates are
\(\gamma_a=-\ell_a\).  The covariance derivative is a congruence, so its
asymptotic growth rates are the allowed pair sums \(\ell_a+\ell_b\), not the
individual one-leg values.  Every reported ``tangent gap'' must state which
of these two spectra is used.  A burn-in interval evolves the base covariance
without accumulating \(R_t\).  The tangent frame is then initialized, evolved
through a declared alignment interval, and retained while its QR log
accumulator is reset before the observation window begins.

\subsection{Maximally mixed input and endpoint identities}

The maximally mixed covariance is \(C_0=\Id/2\), or \(G_0=0\).  For a
regular even word with one-particle transfer product \(M_D\), normalized
state evolution gives the finite-time identity
\begin{equation}
    \widetilde J_D
    =2M_D\bigl(\Id+M_D^\dagger M_D\bigr)^{-1}.
    \label{eq:app-maxmix-regular-bridge}
\end{equation}
Thus a transfer singular value \(s=e^x\) maps to
\begin{equation}
    j=\frac{2s}{1+s^2}=\operatorname{sech}x.
    \label{eq:app-radial-tangent-map}
\end{equation}
This relation identifies \(s\) and \(s^{-1}\); the tangent response retains
the distance from \(s=1\), not its radial orientation.  Equation
\eqref{eq:app-maxmix-regular-bridge} must not be applied letter by letter to
a word containing exact projectors or gain/loss.

There is, however, an exact endpoint statement for an arbitrary nonzero
Gaussian word.  Let its normalized Choi occupied plane have a reduced frame
\begin{equation}
    F_K=\begin{pmatrix}A_K\\ B_K\end{pmatrix},
    \qquad \mathcal G_K=F_K^\dagger F_K.
\end{equation}
Then, up to a phase that drops out of the congruence,
\begin{equation}
    \boxed{
    C_D=B_K\mathcal G_K^{-1}B_K^\dagger,
    \qquad
    \widetilde J_D=2A_K\mathcal G_K^{-1}B_K^\dagger,}
    \label{eq:app-maxmix-endpoint-frame}
\end{equation}
and purity of the Choi projector implies
\begin{equation}
    \boxed{
    \widetilde J_D^\dagger\widetilde J_D
    =4C_D(\Id-C_D).}
    \label{eq:app-endpoint-cocycle-occupation}
\end{equation}
The cocycle kernel locates every constrained output direction, while the
endpoint covariance labels that direction occupied or empty.  This is the
precise maximally mixed relation used in the singular theory.

\subsection{Regulated rank-one Choi-covariance composition}
\label{app:regulated-choi-composition}

The rank loss in Eq.~\eqref{eq:app-half-jacobian-null} may also be handled
directly at the level of pure Gaussian Choi covariances.  Write the incoming
two-leg covariance and the local rank-one relation as
\begin{equation}
    \Sigma=
    \begin{pmatrix}
      \Sigma_{LL}&\Sigma_{LR}\\
      \Sigma_{RL}&\Sigma_{RR}
    \end{pmatrix},
    \qquad
    \mathcal K(\eta_1,\eta_2)=
    \begin{pmatrix}
      \eta_1P&Q\\ Q&\eta_2P
    \end{pmatrix},
    \label{eq:app-local-choi-relation}
\end{equation}
where \(P=\lvert\chi\rangle\langle\chi\rvert\), \(Q=\Id-P\), and
\(\eta_1,\eta_2\in\{+1,-1\}\) encode the selected and reset occupations.
Let
\begin{equation}
    R=\Sigma_{RR},
    \qquad
    d=\Tr(PR)-\eta_1 .
    \label{eq:app-choi-rank-one-denominator}
\end{equation}

\paragraph*{Proposition (finite regulated contraction).}
Replace the singular incoming block by the projector-adapted regulator
\begin{equation}
    \mathcal K_{LL}^{(\epsilon)}=\eta_1P+\epsilon Q.
    \label{eq:app-projector-adapted-regulator}
\end{equation}
If \(d\neq0\), the \(\epsilon\to0\) contraction exists and is
\begin{equation}
\boxed{
\begin{aligned}
\Sigma'_{LL}
&=\Sigma_{LL}-\frac{\Sigma_{LR}P\Sigma_{RL}}{d},\\
\Sigma'_{LR}
&=\Sigma_{LR}Q-\frac{\Sigma_{LR}PRQ}{d},\\
\Sigma'_{RL}
&=Q\Sigma_{RL}-\frac{QRP\Sigma_{RL}}{d},\\
\Sigma'_{RR}
&=\eta_2P+QRQ-\frac{QRPRQ}{d}.
\end{aligned}}
\label{eq:app-regulated-choi-update}
\end{equation}
The regulator acts only on \(\operatorname{Ran}Q\) and retains finite terms generated by
the divergent complementary inverse.  Equation~\eqref{eq:app-regulated-choi-update}
therefore cannot be obtained by blindly replacing the singular block by its
Moore--Penrose pseudoinverse.  When \(d=0\), the displayed finite limit does
not exist; this is a zero-overlap or singular selected branch and requires a
separate rescaled-limit analysis.

If \(\Sigma=\Sigma^\dagger\) and \(\Sigma^2=\Id\), then
Eq.~\eqref{eq:app-regulated-choi-update} obeys
\begin{equation}
    \Sigma'=\Sigma'^\dagger,
    \qquad
    (\Sigma')^2=\Id .
    \label{eq:app-regulated-choi-involution}
\end{equation}
Indeed, choose a basis adapted to \(Q\oplus P\) and reorder the covariance as
\((L\oplus Q)\oplus P\).  Then
\begin{equation}
    \widetilde\Sigma=
    \begin{pmatrix}X&y\\y^\dagger&r\end{pmatrix},
    \qquad
    \widetilde\Sigma'=
    \begin{pmatrix}T&0\\0&\eta_2\end{pmatrix},
    \qquad
    T=X-\frac{yy^\dagger}{r-\eta_1},
    \label{eq:app-choi-involution-proof-blocks}
\end{equation}
with \(r=\Tr(PR)\).  The original involution gives
\begin{equation}
    X^2+yy^\dagger=\Id,
    \qquad Xy+ry=0,
    \qquad y^\dagger y+r^2=1 .
\end{equation}
Using these identities and \(\eta_1^2=1\),
\begin{align}
    T^2
    &=\Id+\left[-1+\frac{2r}{d}
      +\frac{1-r^2}{d^2}\right]yy^\dagger
      =\Id,
    \label{eq:app-choi-Schur-involution}
\end{align}
because \(-d^2+2rd+1-r^2=1-\eta_1^2=0\); together with
\(\eta_2^2=1\), this proves Eq.~\eqref{eq:app-regulated-choi-involution}.

For completeness, the same pure Choi covariance supplies two complementary
transfer charts without assigning an inverse to a projected direction.  For
equal-dimensional legs write
\begin{equation}
    \Sigma=\begin{pmatrix}A&B\\B^\dagger&D\end{pmatrix}.
\end{equation}
On their respective finite sectors, the particle and hole charts are
\begin{equation}
    \mathcal T_p=-(\Id+A)^{-1}B,
    \qquad
    \mathcal T_h=(\Id-A^{\mathsf T})^{-1}B^* .
    \label{eq:app-particle-hole-transfer-charts}
\end{equation}
The particle chart has an endpoint pole at \(A=-\Id\), whereas the hole chart
has the complementary pole at \(A=+\Id\).  For every nonzero finite coupled
singular value,
\begin{equation}
    s_j(\mathcal T_h)=s_j(\mathcal T_p)^{-1}.
    \label{eq:app-particle-hole-reciprocity}
\end{equation}
Thus a projected endpoint is covered by the complementary chart together
with its occupied or empty cap.  Equations~\eqref{eq:app-particle-hole-transfer-charts}
and \eqref{eq:app-particle-hole-reciprocity} are only the chart transition;
the quotient-core and radial-response data are given below and are not
reconstructed here.

\subsection{Caps, quotient core, and trace anchor}
\label{app:cap-core-anchor}

A singular Gaussian word is more naturally a Choi relation than a square
transfer matrix.  For the reduced frame above, define the occupied plane
\begin{equation}
    W_K=\{(A_Kv,B_Kv):v\in\mathbb C^{r_K}\}
    \subset\mathcal H_L\oplus\mathcal H_R.
\end{equation}
Its four intrinsic one-leg cap spaces are
\begin{equation}
    \begin{aligned}
      \mathcal O_R&=B_K\ker A_K,
      &\mathcal E_R&=(\operatorname{Ran}B_K)^\perp,\\
      \mathcal O_L&=A_K\ker B_K,
      &\mathcal E_L&=(\operatorname{Ran}A_K)^\perp.
    \end{aligned}
    \label{eq:app-four-caps}
\end{equation}
They are respectively the exact one- and zero-eigenspaces of the diagonal
Choi covariance blocks.  After these forced directions are removed, the
remaining relation is the graph of an invertible quotient map
\begin{equation}
    R_{\rm core}:\mathcal H_{\rm in}^{f}
    \longrightarrow\mathcal H_{\rm out}^{f},
    \qquad
    T_{\rm core}=R_{\rm core}^\dagger.
    \label{eq:app-quotient-core}
\end{equation}
The final datum is the trace anchor
\begin{equation}
    \nu_K=2^{-N}\Tr(K^\dagger K),
    \label{eq:app-trace-anchor}
\end{equation}
which is absent from the normalized Choi plane but indispensable for absolute
Born probabilities.  Caps plus \(R_{\rm core}\) determine the normalized
endpoint geometry; they do not determine \(\nu_K\).

The centered Choi covariance
\begin{equation}
    \mathcal S_K=2\mathcal C_K-\Id,
    \qquad
    \mathcal S_K^\dagger=\mathcal S_K,
    \qquad
    \mathcal S_K^2=\Id,
    \label{eq:app-auxiliary-scatterer}
\end{equation}
may be read as an auxiliary scattering matrix.  On a free singular channel
of \(T_{\rm core}\) with \(s=e^x\), it reduces to
\begin{equation}
    \mathcal S(x)=
    \begin{pmatrix}
      -\tanh x&\operatorname{sech}x\\
      \operatorname{sech}x&\tanh x
    \end{pmatrix}.
    \label{eq:app-radial-scatterer}
\end{equation}
The transmission amplitude is the maximally mixed tangent singular value in
Eq.~\eqref{eq:app-radial-tangent-map}.  Occupied and empty caps are perfectly
reflecting channels with opposite orientation.

This scattering interpretation is auxiliary: \(\mathcal S_K\) is a
coordinate representation of a Choi boundary relation, not the physical
time-evolution operator of the fermions.  The actual reservoir operation is
the local fSWAP in Eq.~\eqref{eq:app-fswap-unitary}.  Away from maximally
mixed input the tangent factors are dressed by the evolving trajectory and
cannot be identified with fixed transmission blocks.  Likewise, a
topological invariant of an auxiliary scatterer is meaningful only after
its locality, symmetry, gap, and stability have been established; it is not
implied merely by writing Eq.~\eqref{eq:app-auxiliary-scatterer}.

\appendixplaceholderfigure{fig:app-cap-core}{Singular Choi geometry and normalized response}{
Panel (a): the four occupied/empty cap spaces and the finite transmitted
quotient core.  Panel (b): radial free-channel block, showing that
\(s\leftrightarrow s^{-1}\) reverses reflection orientation but leaves the
tangent transmission \(2s/(1+s^2)\) unchanged.  Panel (c): the chronological
state-dressed tangent product and QR accumulation.  The caption should state
explicitly that the Choi scatterer is auxiliary and that the physical
source/drain primitive is fSWAP with a fresh local ancilla.}


\section{Entangling capacity of the local adaptive circuit}
\label{app:adaptive-depth-bound}

We record the preparation bound of Lu, Lessa, Kim, and
Hsieh~\cite{LuLessaKimHsieh2022} in a form adapted to the present singular
instrument.  Consider a pure trajectory and a bipartition
\(\mathcal H_A\widehat\otimes\mathcal H_{\bar A}\).  Let \(\chi_A\) be its
Schmidt rank and \(S_0(A)=\log\chi_A\) its Hartley entropy.  A branch
operator supported entirely within \(A\) or \(\bar A\) cannot increase
\(\chi_A\).  If a crossing operator has operator-Schmidt decomposition
\begin{equation}
    O_j=\sum_{\alpha=1}^{r_j}
    \sigma_{j\alpha}O^A_{j\alpha}\otimes O^{\bar A}_{j\alpha},
    \label{eq:app-operator-schmidt-decomposition}
\end{equation}
then it can multiply the state Schmidt rank by at most \(r_j\).  Consequently
every nonzero adaptive branch obeys
\begin{equation}
    \boxed{
    S_0^{\rm out}(A)-S_0^{\rm in}(A)
    \le\sum_{j\ {\rm crossing}\ \partial A}\log r_j.}
    \label{eq:app-exact-cut-capacity}
\end{equation}
Normalization by the branch probability is a scalar and does not affect this
bound.  Neither invertibility nor a lower bound on the Born probability is
assumed.

For a depth-\(D\) circuit whose layers consist of nonoverlapping,
bounded-support gates, the number of crossing gates per layer is
\(O(|\partial A|)\), while \(r_j=O(1)\).  Equation
\eqref{eq:app-exact-cut-capacity} then gives
\begin{equation}
    S_0^{\rm out}(A)-S_0^{\rm in}(A)
    \le O(D|\partial A|).
    \label{eq:app-Lu-area-depth-bound}
\end{equation}
For a pure product input, \(S_0^{\rm in}=0\), and every positive-index
R\'enyi entropy satisfies
\begin{equation}
    S_n^{\rm out}(A)\le S_0^{\rm out}(A)
    \le O(D|\partial A|),\qquad n>0.
    \label{eq:app-Renyi-depth-bound}
\end{equation}
The product-input qualification is essential.  A Haar-random Slater
determinant is already spatially entangled, while the maximally mixed state
is not a pure state to which Eq.~\eqref{eq:app-Renyi-depth-bound} directly
applies.  The latter may be treated by adjoining a colocated purifier, but
the cut and the retained reference modes must then be specified.

\subsection{Operator-Schmidt refinement for a corrected orbital}

The rank-one fermion instrument permits a sharper count.  Across a spatial
cut, decompose \(a_\chi\) in the graded tensor product as
\begin{equation}
    a_\chi=a_A\otimes\Id+(-1)^{F_A}\otimes a_{\bar A}.
    \label{eq:app-fermion-cut-decomposition}
\end{equation}
It follows that
\begin{equation}
    r(a_\chi),\,r(a_\chi^\dagger)\le2,
    \qquad
    r(n_\chi),\,r(\Id-n_\chi)\le4.
    \label{eq:app-branch-operator-schmidt-ranks}
\end{equation}
Let \(N_=^\partial(\xi)\) count crossing measurements whose result already
equals the target, and \(N_{\ne}^\partial(\xi)\) count crossing measurements
that are corrected by gain or loss.  Equations
\eqref{eq:app-perfect-correction-kraus-table} and
\eqref{eq:app-exact-cut-capacity} give the branchwise bound
\begin{equation}
    \boxed{
    S_0^\xi(A,T)-S_0(A,0)
    \le 2\log2\,N_=^\partial(\xi)
      +\log2\,N_{\ne}^\partial(\xi).}
    \label{eq:app-record-resolved-cut-bound}
\end{equation}
This is a useful implementation check because every quantity on the
right-hand side is available from the same trajectory-resolved record.

\subsection{Full-width strip and the two-dimensional resource}

The entropy used to diagnose the wall is computed for the full-width strip
\begin{equation}
    A_\ell=\{(x,y):0\le x<N_x,\ 0\le y<\ell\}.
    \label{eq:app-full-width-strip}
\end{equation}
On a torus, \(\partial A_\ell\) contains two transverse cuts of length
\(N_x\).  Thus \(|\partial A_\ell|=O(N_x)\), not \(O(1)\).  A subsequent
restriction of the entanglement contour to a window in \(x\) does not change
the subsystem and is not itself the entropy of a wall-only region.

Suppose four rank-one controller orbitals are applied at every active center,
their square support has radius \(r\), and the active slab has transverse
width \(W\).  Away from the short-interval regime in which the two cuts
overlap, at most \(2r\) center rows cross each cut.  Hence one complete sweep
contains at most
\begin{equation}
    N_\partial(T)\le16rWT
    \label{eq:app-crossing-event-count}
\end{equation}
crossing letters after \(T\) sweeps.  From a pure product input,
\begin{equation}
    S_n(A_\ell,T)\le32rWT\log2.
    \label{eq:app-strip-capacity-bound}
\end{equation}
For two oppositely oriented complex-fermion walls, the proposed CFT form is
\begin{equation}
    S_1(A_\ell)
    =b(N_x)+\frac{c_L+c_R}{6}
      \log\!\left[
      \frac{N_y}{\pi a}\sin\!\left(\frac{\pi\ell}{N_y}\right)
      \right]+\cdots ,
    \label{eq:app-two-wall-CFT-form}
\end{equation}
with \(c_L=c_R=1\) in the ideal benchmark.  If this form is established
from a pure product input, Eqs.~\eqref{eq:app-strip-capacity-bound} and
\eqref{eq:app-two-wall-CFT-form} imply the asymptotic resource tradeoff
\begin{equation}
    WT=\Omega(\log N_y).
    \label{eq:app-width-time-tradeoff}
\end{equation}
The intrinsic one-dimensional result \(D=\Omega(\log N_y)\) is recovered
when the transverse resource is fixed.  In a two-dimensional aspect-ratio
limit, the growing transverse boundary supplies entangling capacity.

\subsection{Raster sweeps and finite-color schedules}

The variable \(T\) in Eq.~\eqref{eq:app-crossing-event-count} counts complete
controller sweeps; it is not automatically the circuit depth \(D\) in
Eq.~\eqref{eq:app-Lu-area-depth-bound}.  A raster sweep applies overlapping
OW projectors in a prescribed serial order.  Because neighboring projectors
are generically noncommuting, reproducing that precise word has an extensive
causal depth.

For \(n_{\rm shell}=1\), let \(u=(x,y)\) denote an update center and write its
cell support as
\begin{equation}
    \mathcal S_u
    =\bigl\{u+(\delta_x,\delta_y):
      \delta_x,\delta_y\in\{-1,0,1\}\bigr\},
    \label{eq:app-nshell-one-support}
\end{equation}
with both coordinates understood periodically.  The support-conflict graph
\(\mathcal G_{\rm conf}=(\mathcal V,\mathcal E)\) has one vertex for each
active center and an edge \((u,v)\in\mathcal E\) precisely when
\(\mathcal S_u\cap\mathcal S_v\ne\varnothing\).  In terms of the cyclic
distance
\begin{equation}
    d_L(a,b)=\min_{m\in\mathbb Z}|a-b+mL|,
\end{equation}
two distinct centers conflict if and only if
\begin{equation}
    d_{N_x}(x,x')\le2
    \quad\hbox{and}\quad
    d_{N_y}(y,y')\le2.
    \label{eq:app-conflict-condition}
\end{equation}
The degree is therefore at most \(24\), so a finite coloring with at most
\(25\) colors exists even without using the translational structure.

On a torus for which \(q_x\mid N_x\), \(q_y\mid N_y\), and
\(q_x,q_y\ge3\), the residue map
\begin{equation}
    c(x,y)=(x\bmod q_x,y\bmod q_y)
    \label{eq:app-residue-coloring}
\end{equation}
is a proper \(q_xq_y\)-coloring.  Indeed, two distinct centers of the same
color differ by a nonzero multiple of \(q_x\) or \(q_y\) in at least one
coordinate.  Their cyclic separation in that coordinate is at least three,
so Eq.~\eqref{eq:app-conflict-condition} fails and the supports are disjoint.
Thus a commensurate \(3\times3\) coloring gives nine colors.  For the proposed
sequence
\begin{equation}
    N_x=20,
    \qquad N_y\in\{24,32,40,48,64,80\},
\end{equation}
the conservative choice \(q_x=q_y=4\) wraps periodically and gives
\begin{equation}
    \chi_{\rm conf}\le16.
    \label{eq:app-sixteen-color-bound}
\end{equation}

All centers of a fixed color can be updated in parallel.  If the four
controller orbitals are kept as separate instrument blocks and one
occupation-measurement--feedforward--fSWAP block has bounded physical depth
\(d_{\rm inst}=O(1)\), a complete colored sweep obeys
\begin{equation}
    d_{\rm cyc}
    \le4\chi_{\rm conf}d_{\rm inst}
    \le64d_{\rm inst}.
    \label{eq:app-colored-cycle-depth}
\end{equation}
The bound is independent of \(N_x\) and \(N_y\).  Consequently
\(D=d_{\rm cyc}T_{\rm cyc}\), and the two-wall entropy
\(S_1\sim\frac13\log N_y\), together with
Eq.~\eqref{eq:app-Lu-area-depth-bound}, implies
\begin{equation}
    T_{\rm cyc}
    \gtrsim
    \frac{\log N_y}{3\kappa d_{\rm cyc}N_x}
    \ge
    \frac{\log N_y}{192\kappa d_{\rm inst}N_x},
    \label{eq:app-colored-cycle-Lu-bound}
\end{equation}
where \(\kappa\) absorbs the local operator-Schmidt constants and the two
components of \(\partial A_\ell\).  At fixed \(N_x\), this rules out a
sublogarithmic number of colored cycles for preparation from a product state,
but its unknown prefactor is far too weak to prescribe a practical burn-in.
The constant bound in Eq.~\eqref{eq:app-colored-cycle-depth} is per colored
sweep only.  It does not make the total depth \(D\) size independent, and the
BPJ result on rapid bulk preparation cannot be transferred without proof to
the additional logarithmically entangled wall sector.

This implementation is not a reinterpretation of raster data.  If \(u\) and
\(v\) overlap, their conditional instruments generally satisfy
\begin{equation}
    \mathcal M_u\circ\mathcal M_v
    \ne
    \mathcal M_v\circ\mathcal M_u.
    \label{eq:app-instrument-noncommutativity}
\end{equation}
Changing from the raster word to color blocks therefore defines a different
microscopic circuit, even though both use the same local controller.  An
orientation-unbiased implementation may draw a fresh permutation of the
color blocks each cycle, updating all sites within one block in parallel;
the color order and its random seed must be recorded.  Raster and
reverse-raster schedules remain useful artifact controls.  The colored
protocol can replace the serial protocol in production only after topology,
entanglement and charge coefficients, chirality, and the tangent slow sector
agree within predeclared finite-size confidence intervals.

\subsection{Finite wall separation and hybridization}

On a periodic transverse direction there are two oppositely oriented walls.
If their minimum separation is \(w_{\min}\), exponentially localized edge
modes generically acquire a hybridization mass
\begin{equation}
    m(w_{\min})\simeq m_0e^{-w_{\min}/\xi_\perp}.
    \label{eq:app-wall-hybridization-mass}
\end{equation}
The corresponding low-energy two-wall Hamiltonian has the Dirac form
\begin{equation}
    h_{\rm edge}(k_y)=v k_y\tau_z+m(w_{\min})\tau_x.
    \label{eq:app-two-wall-Dirac}
\end{equation}
Critical finite-size scaling on a circumference \(N_y\) requires the
hybridization scale to remain below the level spacing,
\begin{equation}
    m(w_{\min})\ll\frac{2\pi v}{N_y}.
\end{equation}
Therefore an effectively decoupled-wall sequence should satisfy
\begin{equation}
    \boxed{
    w_{\min}\gtrsim\xi_\perp
    \log\!\left(\frac{m_0N_y}{2\pi v}\right).}
    \label{eq:app-wall-width-condition}
\end{equation}
This physical width requirement mirrors
Eq.~\eqref{eq:app-width-time-tradeoff}.  A controlled test should extract
\(\xi_\perp\) independently from the transverse wall profile and seek a
collapse in the variable \(N_ye^{-w_{\min}/\xi_\perp}\).

\appendixplaceholderfigure{fig:app-resource-bound}{Where the preparation resource resides}{
Panel (a): a full-width strip with its two transverse entanglement cuts and
two domain walls.  Panel (b): crossing branch operators counted in
Eq.~\eqref{eq:app-record-resolved-cut-bound}.  Panel (c): raster versus
finite-color scheduling.  Panel (d): the two-wall hybridization crossover
controlled by \(N_y e^{-w_{\min}/\xi_\perp}\).  All plotted curves and fit
parameters are placeholders until the corresponding campaigns are complete.}


\section{Exact-Hamiltonian benchmark and boundary twists}
\label{app:exact-H-benchmark}

The equilibrium reference is the two-band class-A lattice Hamiltonian
\begin{equation}
    h(\mathbf k;\alpha)
    =\sin k_x\,\sigma_x+\sin k_y\,\sigma_y
     +\bigl(\alpha-\cos k_x-\cos k_y\bigr)\sigma_z.
    \label{eq:app-QWZ-Hamiltonian}
\end{equation}
A spatial profile \(\alpha(x)\) is Fourier transformed only in the
translation-invariant direction, or assembled directly as a real-space
Hermitian matrix \(H_{\rm ex}\) on \(2N_xN_y\) orbitals.  The exact benchmark
contains no Born sampling and no feedback.

At half filling, diagonalize
\begin{equation}
    H_{\rm ex}V=V\,\operatorname{diag}(\varepsilon_1,\ldots,\varepsilon_{2N_xN_y}),
    \qquad
    \varepsilon_1\le\cdots\le\varepsilon_{2N_xN_y},
\end{equation}
and occupy exactly \(N_xN_y\) orbitals,
\begin{align}
    P_{\rm ex}&=V_{\rm occ}V_{\rm occ}^\dagger,
    &C_{\rm ex}&=P_{\rm ex}^*,\nonumber\\
    G_{\rm ex}&=2C_{\rm ex}-\Id,
    &V_{\rm occ}&=(v_1,\ldots,v_{N_xN_y}).
    \label{eq:app-exact-half-filled-projector}
\end{align}
Here \(P_{\rm ex}\) is the ket-space spectral projector, whereas
\(C_{\rm ex}\) is the matrix stored with the convention of
Eq.~\eqref{eq:app-C-G-convention}.  Topological estimators act on
\(P_{\rm ex}=C_{\rm ex}^{\mathsf T}\); covariance observables act on
\(C_{\rm ex}\).
Selecting states merely by the sign of \(\varepsilon\) is ambiguous when a
finite-size wall level sits at zero.  Fixed particle number is the primary
half-filling convention.  Any residual degeneracy at the Fermi level must be
resolved by a stated symmetry, an infinitesimal twist, or explicit
subspace tracking.

\subsection{Physical and modular current densities}
\label{app:physical-modular-current}

For a general number-conserving quadratic Hamiltonian, introduce the orbital
spinor \(\Psi_{\mathbf r}\) in cell \(\mathbf r\) and write
\begin{equation}
    H=\sum_{\mathbf r,\mathbf r'}
      \Psi_{\mathbf r}^\dagger
      h_{\mathbf r\mathbf r'}\Psi_{\mathbf r'},
    \qquad
    n_{\mathbf r}=\Psi_{\mathbf r}^\dagger\Psi_{\mathbf r}.
    \label{eq:app-general-quadratic-H}
\end{equation}
Hermiticity requires
\(h_{\mathbf r'\mathbf r}=h_{\mathbf r\mathbf r'}^\dagger\).  The
Heisenberg equation gives the lattice continuity equation
\begin{equation}
    \dot n_{\mathbf r}
    +\sum_{\mathbf r'}J_{\mathbf r\to\mathbf r'}=0,
    \label{eq:app-lattice-continuity}
\end{equation}
with the positively oriented, outward bond-current operator
\begin{equation}
    J_{\mathbf r\to\mathbf r'}
    =\ii\Psi_{\mathbf r}^\dagger
       h_{\mathbf r\mathbf r'}\Psi_{\mathbf r'}
     -\ii\Psi_{\mathbf r'}^\dagger
       h_{\mathbf r'\mathbf r}\Psi_{\mathbf r}.
    \label{eq:app-bond-current-operator}
\end{equation}
Using the covariance convention of Eq.~\eqref{eq:app-C-G-convention}, its
Gaussian expectation value is
\begin{equation}
    \boxed{
    j_{\mathbf r\to\mathbf r'}
    =2\,\operatorname{Im}\Tr\!\left[
      h_{\mathbf r\mathbf r'}C_{\mathbf r'\mathbf r}\right],}
    \qquad
    j_{\mathbf r\to\mathbf r'}=-j_{\mathbf r'\to\mathbf r}.
    \label{eq:app-bond-current-expectation}
\end{equation}
For the exact QWZ benchmark of Eq.~\eqref{eq:app-QWZ-Hamiltonian}, the two
positively oriented nearest-neighbor currents are therefore
\begin{align}
    j_x(\mathbf r)
    &=2\,\operatorname{Im}\Tr\!\left[
      h_{\mathbf r,\mathbf r+\widehat x}
      C_{\mathbf r+\widehat x,\mathbf r}\right],\nonumber\\
    j_y(\mathbf r)
    &=2\,\operatorname{Im}\Tr\!\left[
      h_{\mathbf r,\mathbf r+\widehat y}
      C_{\mathbf r+\widehat y,\mathbf r}\right].
    \label{eq:app-QWZ-link-currents}
\end{align}

The same construction defines a modular current for a subsystem \(A\).  On
the entanglement-active support \(0<C_A<\Id_A\), the one-body modular
Hamiltonian is
\begin{align}
    h_A^{\rm mod}
    &=\log\!\left[(\Id_A-C_A)C_A^{-1}\right]\nonumber\\
    &=-2\arctanh(G_A),
    &G_A&=2C_A-\Id_A .
    \label{eq:app-modular-Hamiltonian}
\end{align}
Replacing \(h_{\mathbf r\mathbf r'}\) in
Eq.~\eqref{eq:app-bond-current-expectation} by the corresponding blocks of
\(h_A^{\rm mod}\), while retaining \(C_A\), gives the modular-current map.
Exact eigenvalues \(0\) or \(1\) lie outside this finite chart.  Numerically
they are either removed by restricting to the entanglement-active support or
clipped to \([\epsilon,1-\epsilon]\); any spatial current claim must be stable
as \(\epsilon\) is reduced.

Equations~\eqref{eq:app-lattice-continuity}--\eqref{eq:app-QWZ-link-currents}
are source free for exact-Hamiltonian and modular evolution.  The adaptive
system alone exchanges charge with measured and fSWAP ancillas, so a
circuit-time balance law contains a local source term,
\begin{equation}
    \Delta n_{\mathbf r}
    +(\nabla_{\rm lat}\!\cdot j)_{\mathbf r}
    =s_{\mathbf r}^{\rm anc}.
    \label{eq:app-adaptive-current-source}
\end{equation}
Consequently a source-free physical current is not inferred from the
measurement/feedforward occupancy change; charge conservation holds instead
on the enlarged system--ancilla Hilbert space.

\subsection{Twisted boundary condition}

A twist \(\phi_y\) is inserted into the same Hamiltonian by multiplying the
forward hopping across one \(y\)-seam by \(e^{i\phi_y}\) and the reverse
hopping by \(e^{-i\phi_y}\).  A uniform gauge distributes the phase as
\(e^{i\phi_y/N_y}\) over all forward \(y\)-links and is unitarily
equivalent.  The spatial profile, disorder realization, and every other
parameter are held fixed as \(\phi_y\) varies.

Let \(V_j=V_{\rm occ}(\phi_j)\) be an orthonormal occupied frame.  Neighboring
frames are aligned through the singular-value decomposition
\begin{equation}
    M_j=V_j^\dagger V_{j+1}=L_j\Sigma_jR_j^\dagger,
    \qquad
    \mathcal U_j=L_jR_j^\dagger .
    \label{eq:app-overlap-polar-link}
\end{equation}
Multiplying \(V_{j+1}\) by \(R_jL_j^\dagger\) makes the overlap with \(V_j\)
positive and fixes a discrete parallel-transport gauge.  The Wilson
holonomy and its many-body phase are
\begin{equation}
    \mathcal W_y=\mathcal U_0\mathcal U_1\cdots\mathcal U_{N_\phi-1},
    \qquad
    \gamma_y=\arg\det\mathcal W_y.
    \label{eq:app-exact-Wilson-holonomy}
\end{equation}
This polar alignment is a gauge choice for one continuous family.  It cannot
turn independently sampled Hamiltonians or independently sampled circuit
records into such a family.

If the half-filled gap remains open, \(P_{\rm ex}(\phi_y)\) defines a smooth
ground-state bundle.  At an exact spectral crossing, there are two distinct
questions: the instantaneous half-filled ground-state projector may switch
discontinuously, whereas a maximally overlapping tracked subspace follows a
spectral-flow branch and need not remain the instantaneous ground state.
Both are useful, but they must be labeled separately.

For a two-twist grid \(\boldsymbol\phi=(\phi_x,\phi_y)\), a gauge-invariant
lattice Chern number is obtained from determinant links
\begin{equation}
    U_\mu(\boldsymbol\phi)
    =\frac{\det[V^\dagger(\boldsymbol\phi)
      V(\boldsymbol\phi+\widehat\mu)]}
    {|\det[V^\dagger(\boldsymbol\phi)
      V(\boldsymbol\phi+\widehat\mu)]|},
    \label{eq:app-FHS-link}
\end{equation}
and plaquette curvature
\begin{align}
    F_{xy}(\boldsymbol\phi)
    &=\operatorname{Arg}\!\left[
      U_x(\boldsymbol\phi)
      U_y(\boldsymbol\phi+\widehat x)
      U_x(\boldsymbol\phi+\widehat y)^{-1}
      U_y(\boldsymbol\phi)^{-1}\right],\nonumber\\
    C_{\rm tw}&=\frac{1}{2\pi}\sum_{\boldsymbol\phi}F_{xy}.
    \label{eq:app-twist-Chern-number}
\end{align}
The mesh must be refined until the admissibility condition and the integer
are stable~\cite{FukuiHatsugaiSuzuki2005}.

The exact benchmark comparison should report separately: the half-filled
spectrum and wall velocity, the transverse localization length, the strip
entropy and contour, the local or twisted-boundary Chern diagnostic, and the
finite-width hybridization gap.  Agreement of a circuit observable with one
benchmark does not imply equality of the two states.

\appendixplaceholderfigure{fig:app-exact-benchmark}{Exact Hamiltonian, half filling, and twist tracking}{
Panel (a): the real-space domain-wall profile and exact low-energy spectrum.
Panel (b): occupied-subspace tracking under a boundary twist, including the
distinction between instantaneous ground-state filling and a spectral-flow
branch.  Panel (c): convergence of the twist-mesh Chern estimator.  Panel
(d): proposed circuit-versus-exact comparisons.  Numerical values are left
blank until established.}


\section{Trajectory estimators and averaging conventions}
\label{app:estimators}

Every stochastic estimator below is evaluated on independently Born-sampled
records.  Equal sample weights therefore estimate the physical Born ensemble;
no second factor of \(p_\xi\) is inserted.  We use an overbar for this Monte
Carlo average,
\begin{equation}
    \overline{X_\xi}=\frac{1}{S}\sum_{\xi=1}^{S}X_\xi.
    \label{eq:app-Born-sample-average}
\end{equation}
Nonlinear functions are evaluated trajectory by trajectory unless an
explicitly annealed quantity is stated.

\subsection{Entropy and entanglement contour}

For a Gaussian trajectory and a spatial subsystem \(A\), let
\(C_{\xi,A}\) be the restricted two-point matrix.  The Gaussian von Neumann
entropy is
\begin{align}
    S_\xi(A)
    &=-\operatorname{tr}[C_{\xi,A}\log C_{\xi,A}]\nonumber\\
    &\quad-\operatorname{tr}[(\Id_A-C_{\xi,A})
      \log(\Id_A-C_{\xi,A})].
    \label{eq:app-Gaussian-entropy}
\end{align}
If
\(C_{\xi,A}u_{\xi,\alpha}=\nu_{\xi,\alpha}u_{\xi,\alpha}\), define
\(h(\nu)=-\nu\log\nu-(1-\nu)\log(1-\nu)\).  A site-resolved contour is
\begin{equation}
    s_{\xi,A}(i)
    =\sum_\alpha |u_{\xi,\alpha}(i)|^2h(\nu_{\xi,\alpha})
    =\bigl[h(C_{\xi,A})\bigr]_{ii}.
    \label{eq:app-entanglement-contour}
\end{equation}
Internal orbitals in one unit cell are summed only after
Eq.~\eqref{eq:app-entanglement-contour}.  It follows exactly that
\begin{equation}
    S_\xi(A)=\sum_{i\in A}s_{\xi,A}(i).
\end{equation}
For a mixed trajectory, Eq.~\eqref{eq:app-Gaussian-entropy} is the entropy of
the reduced mixed state, not an entanglement measure.  Purity, or an
appropriate mixed-state entanglement diagnostic, must accompany any such
interpretation.

For the full-width strip \(A_\ell\), a wall-window contour sum is
\begin{equation}
    \mathcal S_{\xi,I_x}(\ell)
    =\sum_{x\in I_x}\sum_{y=0}^{\ell-1}
      \sum_\mu s_{\xi,A_\ell}(x,y,\mu).
    \label{eq:app-windowed-contour}
\end{equation}
Only \(I_x=[0,N_x-1]\) gives the entropy of \(A_\ell\).  A narrow \(I_x\)
is a localization diagnostic for the contour of the same full-width
subsystem; it is not the entropy of a narrow tube.

The primary trajectory estimator is
\begin{equation}
    \overline{S(A_\ell)}
    =\frac1S\sum_\xi S_\xi(A_\ell),
    \label{eq:app-average-trajectory-entropy}
\end{equation}
which generally differs from both \(S[\overline{C_A}]\) and the entropy of
the averaged many-body density matrix.  A proposed two-wall CFT fit uses
\begin{align}
    d_{N_y}(\ell)
    &=\frac{N_y}{\pi}
      \sin\!\left(\frac{\pi\ell}{N_y}\right),\nonumber
    \\
    \overline{S(A_\ell)}
    &=b+\frac{c_L+c_R}{6}\log d_{N_y}(\ell)+\cdots .
    \label{eq:app-entropy-log-chord-fit}
\end{align}
Fit windows, endpoint exclusions, finite-size corrections, and covariance of
the fitted points must be specified.  A coefficient inferred from
\(\mathcal S_{I_x}\) for one wall is supporting contour evidence, not a
standalone entropy measurement.

\subsection{Conditional mutual information and the tripartite reduction}
\label{app:entropy-tmi-relation}

Let \(A,B,C\) be adjacent intervals in the periodic wall direction, with each
region extending across the same full transverse width.  For every trajectory
\(\xi\), define the mutual information between the separated intervals and
the standard tripartite information by
\begin{align}
    I_{2,\xi}(A,C)
    &=S_\xi(A)+S_\xi(C)-S_\xi(A\cup C),\nonumber\\
    I_{3,\xi}(A,B,C)
    &=S_\xi(A)+S_\xi(B)+S_\xi(C)\nonumber\\
    &\quad-S_\xi(A\cup B)-S_\xi(A\cup C)\nonumber\\
    &\quad-S_\xi(B\cup C)+S_\xi(A\cup B\cup C).
    \label{eq:app-I2-I3-definitions}
\end{align}
Their difference eliminates the disconnected-region entropy exactly:
\begin{align}
    \mathcal I_\xi(A:C|B)
    &\equiv I_{2,\xi}(A,C)-I_{3,\xi}(A,B,C)\nonumber\\
    &=S_\xi(A\cup B)+S_\xi(B\cup C)\nonumber\\
    &\quad-S_\xi(B)-S_\xi(A\cup B\cup C)\nonumber\\
    &=I_\xi(A:C|B),
    \label{eq:app-TMI-CMI-identity}
\end{align}
so the quantity called the ``tripartite reduction'' is precisely a
conditional mutual information and is nonnegative by strong subadditivity.

For a common interval family governed by a one-dimensional CFT, the averaged
entropy has the finite-ring form~\cite{CalabreseCardy2004}
\begin{equation}
    \overline{S(X)}
    =\kappa\log d_L(\ell_X)+s_0+\cdots,
    \qquad
    d_L(\ell)=\frac{L}{\pi}\sin\!\left(\frac{\pi\ell}{L}\right),
    \label{eq:app-CFT-entropy-common-family}
\end{equation}
where \(\kappa\) includes the number and orientation-independent content of
the wall branches intersected by the chosen geometry.  Substitution into
Eq.~\eqref{eq:app-TMI-CMI-identity} cancels \(s_0\) and gives
\begin{equation}
    \overline{\mathcal I(A:C|B)}
    =\kappa\log\!\left[
      \frac{d_L(A\cup B)d_L(B\cup C)}
           {d_L(B)d_L(A\cup B\cup C)}\right]+\cdots .
    \label{eq:app-CMI-chord-ratio}
\end{equation}
Define the chord cross ratio
\begin{equation}
    x=\frac{d_L(A)d_L(C)}
            {d_L(A\cup B)d_L(B\cup C)}.
    \label{eq:app-CMI-cross-ratio}
\end{equation}
For adjacent lengths \(\ell_A,\ell_B,\ell_C\), the elementary identity
\begin{equation}
    \sin(a+b)\sin(b+c)-\sin a\sin c
    =\sin b\sin(a+b+c)
\end{equation}
with \(a=\pi\ell_A/L\), \(b=\pi\ell_B/L\), and
\(c=\pi\ell_C/L\) implies
\begin{equation}
    d_L(A\cup B)d_L(B\cup C)-d_L(A)d_L(C)
    =d_L(B)d_L(A\cup B\cup C).
    \label{eq:app-chord-cross-ratio-identity}
\end{equation}
Consequently,
\begin{equation}
    \boxed{
    \overline{I_2-I_3}
    =\kappa\log\!\left(\frac{1}{1-x}\right)+\cdots .}
    \label{eq:app-TMI-cross-ratio-scaling}
\end{equation}

The estimator order is part of this statement.  Each of the four entropies in
Eq.~\eqref{eq:app-TMI-CMI-identity} is computed from the same trajectory
covariance before forming \(\mathcal I_\xi\), and only then is the equal-weight
Born average taken:
\begin{equation}
    \overline{\mathcal I(A:C|B)}
    =\frac1S\sum_{\xi=1}^{S}\mathcal I_\xi(A:C|B).
    \label{eq:app-trajectory-first-CMI-average}
\end{equation}
Using \(S[\overline C]\), or mixing entropies drawn from independently sampled
records, defines a different quantity.  For a genuine single complex chiral
wall with \(c_L+c_R=1\), Eq.~\eqref{eq:app-TMI-cross-ratio-scaling} has
\(\kappa=1/6\).  A full-width strip intersecting the two oppositely oriented
walls has \(\kappa=1/3\).  This factor of two counts wall branches and does not
measure their propagation direction.

The archived exact-domain-wall ground state provides a deterministic
calibration of this estimator at \(N_x\times N_y=20\times40\), with
\(\ell_A=\ell_C=3\).  The fitted slopes are summarized in
Table~\ref{tab:app-exact-TMI-calibration}.
\begin{table}[t]
\caption{Exact-Hamiltonian calibration of
\(I_2-I_3\) against \(\log[1/(1-x)]\).  Absolute fit uncertainties are
displayed explicitly.}
\label{tab:app-exact-TMI-calibration}
\centering
\footnotesize
\begin{tabular}{@{}lcc@{}}
\toprule
Estimator & Fitted slope [CFT] & \(R^2\)\\
\midrule
\makecell[l]{Each wall-window\\contour sum}
& \makecell{\(0.166967668\)\\\(\pm0.000027005\) \([1/6]\)}
& \(0.99999945\)\\
Full-width CMI
& \makecell{\(0.334232464\)\\\(\pm0.000095377\) \([1/3]\)}
& \(0.99999829\)\\
\bottomrule
\end{tabular}
\end{table}
Only the full-width row is a conditional mutual information of spatial
subsystems.  In the first row the subsystem is still defined over the full
transverse width and only the final entanglement-contour sum is restricted to
one wall; those values test localization and the factor of two, not a
standalone one-wall CMI.  Both rows are deterministic exact-Hamiltonian
calibrations.  They do not constitute a production Born-trajectory result for
the adaptive circuit, and Eq.~\eqref{eq:app-TMI-cross-ratio-scaling} is not an
independent central-charge estimator when it is assembled from the same
entropies used in Eq.~\eqref{eq:app-entropy-log-chord-fit}.

\subsection{Quantum charge variance, full counting statistics, and current algebra}

For
\(N_A=\sum_{i\in A}c_i^\dagger c_i\), the charge expectation and the
within-trajectory quantum variance are
\begin{align}
    \langle N_A\rangle_\xi&=\operatorname{tr}C_{\xi,A},\nonumber\\
    F_{\xi,A}^{(q)}
    &=\langle(N_A-\langle N_A\rangle_\xi)^2\rangle_\xi\nonumber\\
    &=\operatorname{tr}\!\left[C_{\xi,A}(\Id_A-C_{\xi,A})\right].
    \label{eq:app-quantum-charge-variance}
\end{align}
The total variance obtained by first drawing a Born trajectory and then
measuring \(N_A\) obeys the law of total variance,
\begin{equation}
    F_A^{\rm total}
    =\overline{F_{\xi,A}^{(q)}}
     +\operatorname{Var}_\xi\!\left(\operatorname{tr}C_{\xi,A}\right).
    \label{eq:app-total-charge-variance}
\end{equation}
The second term is classical record-to-record charge noise.  It must not be
included in a Kac--Moody coefficient intended to characterize quantum charge
fluctuations of the conditional state.

The trajectory full-counting-statistics generator is the determinant
\begin{equation}
    \mathcal Z_{\xi,A}(\lambda)
    =\left\langle e^{i\lambda N_A}\right\rangle_\xi
    =\det\!\left[\Id_A-C_{\xi,A}
      +e^{i\lambda}C_{\xi,A}\right],
    \label{eq:app-FCS-determinant}
\end{equation}
with cumulants
\begin{equation}
    \kappa_{n,\xi}^{(q)}
    =(-i\partial_\lambda)^n
      \log\mathcal Z_{\xi,A}(\lambda)\big|_{\lambda=0}.
    \label{eq:app-FCS-cumulants}
\end{equation}
The Born-ensemble characteristic function
\(\overline{\mathcal Z_{\xi,A}}\) describes the two-stage physical sampling
experiment.  The typical generator
\(\exp[\overline{\log\mathcal Z_{\xi,A}}]\) is a different, quenched
quantity.

With current normalization
\begin{equation}
    \langle J(z)J(0)\rangle=\frac{k}{z^2},
    \label{eq:app-KM-normalization}
\end{equation}
a single chiral branch at level \(k\) contributes
\begin{equation}
    \overline{F_{\xi,A_\ell}^{(q)}}
    =b_F+\frac{k}{2\pi^2}\log d_{N_y}(\ell)+\cdots .
    \label{eq:app-single-wall-charge-fluctuation}
\end{equation}
For two independent opposite walls, the coefficients add:
\begin{equation}
    \overline{F_{\xi,A_\ell}^{(q)}}
    =b_F+\frac{k_L+k_R}{2\pi^2}
      \log d_{N_y}(\ell)+\cdots .
    \label{eq:app-two-wall-charge-fluctuation}
\end{equation}
The ideal complex-fermion benchmark has \(k_L=k_R=1\).  In that free
benchmark the leading terms satisfy
\begin{equation}
    S_{\rm wall}=\frac{\pi^2}{3}F_{\rm wall}^{(q)}+\cdots .
    \label{eq:app-entropy-fluctuation-ratio}
\end{equation}
The paper should treat Eqs.~\eqref{eq:app-two-wall-charge-fluctuation} and
\eqref{eq:app-entropy-fluctuation-ratio} as tests to be performed, rather
than as consequences of the observed entropy coefficient.

\subsection{Correlations and spatial mode densities}

Wick's theorem gives the connected density correlator on one trajectory,
\begin{equation}
    \langle\delta n_i\,\delta n_j\rangle_\xi
    =\delta_{ij}C_{\xi,ii}-|C_{\xi,ij}|^2.
    \label{eq:app-density-correlation}
\end{equation}
For separated cells one may therefore use the positive squared-correlation
diagnostic
\begin{equation}
    \mathcal C_\xi(x,r)
    =\frac{1}{N_y}\sum_y\sum_{\mu,\nu}
      \left|C_{\xi,(x,y,\mu),(x,y+r,\nu)}\right|^2.
    \label{eq:app-squared-correlation}
\end{equation}
The annealed and typical correlations are respectively
\begin{equation}
    \mathcal C_{\rm ann}(x,r)=\overline{\mathcal C_\xi(x,r)},
    \qquad
    \mathcal C_{\rm typ}(x,r)
    =\exp\!\left[\overline{\log\mathcal C_\xi(x,r)}\right].
    \label{eq:app-annealed-typical-correlation}
\end{equation}
They need not share an exponent in a broad stochastic ensemble.  In
particular,
\begin{equation}
    \overline{|C_{\xi,ij}|^2}
    \ne |\overline{C_{\xi,ij}}|^2.
    \label{eq:app-correlation-average-order}
\end{equation}
Every reported log--log fit must state whether it uses
\(\log\overline{\mathcal C_\xi}\) or
\(\overline{\log\mathcal C_\xi}\).

For a normalized tangent or endpoint vector \(v_{\xi,a}\), the physical
component-summed spatial density is
\begin{equation}
    \rho_a(x,y)
    =\frac1S\sum_{\xi=1}^{S}\sum_\mu
      |v_{\xi,a}(x,y,\mu)|^2,
    \qquad
    \sum_{x,y}\rho_a(x,y)=1.
    \label{eq:app-tangent-mode-density}
\end{equation}
Normalizing each internal component \(\mu\) separately before averaging
instead produces the spatial distribution conditioned on that component and
discards its physical relative weight.  Such a visualization, if used, must
be labeled accordingly.

Statistical uncertainty is estimated over independent trajectories, not over
spatial points reused within one trajectory.  Translation averages reduce
variance but do not create additional independent samples.  Bootstrap or
jackknife resampling should preserve the full trajectory as the elementary
block.


\section{Born-ensemble response and fixed-record holonomy}
\label{app:response-holonomy}

\subsection{The score term in a physical response}

Let a perturbation parameter \(u\) modify the microscopic circuit, and let
\(O_\xi(u)\) be an observable on the normalized trajectory selected by
\(\xi\).  The Born-ensemble expectation is
\begin{equation}
    \langle O\rangle_{\rm B}(u)
    =\sum_\xi p_\xi(u)O_\xi(u).
\end{equation}
Differentiating gives the exact likelihood-ratio identity
\begin{equation}
    \boxed{
    \partial_u\langle O\rangle_{\rm B}
    =\overline{\partial_uO_\xi}
     +\overline{O_\xi\,\partial_u\log p_\xi}.}
    \label{eq:app-full-Born-response}
\end{equation}
Since \(\overline{\partial_u\log p_\xi}=0\), the second term may equivalently
be written as a covariance with the score.  The first term is the
fixed-record conditional response; the second accounts for the change in
the probability of drawing that record.  Dropping it does not give the
linear response of the physical Born ensemble.

The score is accumulated chronologically,
\begin{equation}
    \partial_u\log p_\xi
    =\sum_t\partial_u\log
      p_{m_t}\!\left(G_{t-1}(u),\chi_t(u)\right),
    \label{eq:app-record-score-sum}
\end{equation}
where the derivative is total and includes the perturbed earlier covariance.
For the binary measurement of Eq.~\eqref{eq:app-binary-born-probability},
\begin{equation}
    \partial_u\log p_{m_t}
    =\frac{\eta_{m_t}\,\partial_u g_{\chi_t}}
    {1+\eta_{m_t}g_{\chi_t}},
    \label{eq:app-binary-score}
\end{equation}
with
\begin{equation}
    \partial_u g_{\chi}
    =\chi^\dagger(\partial_uG)\chi
     +2\operatorname{Re}
      \left[(\partial_u\chi)^\dagger G\chi\right].
\end{equation}
The tangent recursion supplies \(\partial_uG\), while the second term is
retained when the perturbation changes the measured orbital itself.

For a wall-density response to a localized potential,
\begin{equation}
    \chi_{\rm B}(x,t)
    =\left.\frac{\delta\,\overline{\langle n(x,t)\rangle_\xi}}
    {\delta V(0,0)}\right|_{V=0},
    \label{eq:app-Born-wall-response}
\end{equation}
Eq.~\eqref{eq:app-full-Born-response} fixes the estimator.  Fourier
transformation along the wall may then be used to search for a low-energy
ridge in \(\chi_{\rm B}(k,\omega)\).  A one-sided slope and its reversal
under wall-orientation reversal would diagnose chirality; no quantized
transport coefficient follows from the ridge alone.

\subsection{Fixed-record twist family}

For a circuit record \(\xi\), insert a twist \(\phi\) into the same
finite-support orbitals and fSWAP couplings while holding the sequence of
outcome labels fixed:
\begin{equation}
    K_\xi(\phi)
    =K_{m_D|s_D}(\phi)\cdots K_{m_1|s_1}(\phi).
    \label{eq:app-fixed-record-word}
\end{equation}
The normalized state
\begin{equation}
    \rho_\xi(\phi)
    =\frac{K_\xi(\phi)\rho_0K_\xi(\phi)^\dagger}
    {p_\xi(\phi)}
    \label{eq:app-fixed-record-state}
\end{equation}
defines a continuous family only on intervals where \(p_\xi(\phi)>0\) and
the relevant occupied rank is constant.  A zero of the branch probability
is a genuine singularity of this normalized family.

For a pure Gaussian branch with occupied frame \(V_{\xi,j}\) on a twist
mesh, use the polar links of Eq.~\eqref{eq:app-overlap-polar-link} and form
\begin{equation}
    \mathcal W_\xi
    =\prod_j\operatorname{polar}
      \!\left(V_{\xi,j}^\dagger V_{\xi,j+1}\right),
    \qquad
    \gamma_\xi=\arg\det\mathcal W_\xi.
    \label{eq:app-fixed-record-holonomy}
\end{equation}
For a mixed Gaussian state, the appropriate object is instead a specified
purification, Uhlmann holonomy, or a spectral projector of the modular
Hamiltonian; these choices are not interchangeable.

Independently resampling a Hamiltonian, record, or ground state at every
\(\phi\) does not define Berry transport.  SVD alignment can fix the gauge
within one family, but it cannot manufacture continuity between unrelated
samples.  A single \(0\le\phi_y\le2\pi\) loop measures a holonomy or spectral
flow.  A Chern number requires a two-dimensional twist family, such as
Eq.~\eqref{eq:app-twist-Chern-number}, or an independently justified
real-space invariant.


\section{Conditioned trajectories, mean channel, and record weights}
\label{app:ensemble-distinctions}

For one adaptive step, denote the unnormalized branch map by
\begin{equation}
    \mathcal K_{m|s}(\rho)
    =K_{m|s}\rho K_{m|s}^\dagger.
\end{equation}
There are three distinct objects:
\begin{align}
    \rho_{m|s}
    &=\frac{\mathcal K_{m|s}(\rho)}
      {\Tr\mathcal K_{m|s}(\rho)}
      &&\text{conditioned trajectory},\\
    \mathcal E_s(\rho)
    &=\sum_m\mathcal K_{m|s}(\rho)
      &&\text{unconditional channel},\\
    \rho_{m_\star}
    &=\frac{\Pi_{m_\star}\rho\Pi_{m_\star}}
      {\Tr(\rho\Pi_{m_\star})}
      &&\text{postselected outcome}.
    \label{eq:app-three-ensemble-objects}
\end{align}
The protocol uses the first object with Born-distributed \(m\) and
deterministic feedback.  A no-click non-Hermitian Hamiltonian or forced
measurement retains only a selected branch and is not the same instrument.

The one-step perfect-correction average is the Gaussian replacement channel
in Eq.~\eqref{eq:app-reset-channel-density}.  More general convex mixtures
of Gaussian Kraus maps need not preserve Gaussianity, especially when future
operations depend on coarse-grained records.  In either case the conceptual
distinction remains:
\begin{equation}
    \overline{S(C_\xi)}\ne S(\overline C),
    \qquad
    \overline{|C_{\xi,ij}|^2}\ne|\overline C_{ij}|^2.
    \label{eq:app-mean-trajectory-nonlinearity}
\end{equation}
Even when \(\overline\rho\) happens to be Gaussian, its entropy is the
entropy of the mixed mean state, not the average trajectory entanglement.
The tangent cocycle is the derivative of the normalized conditioned map and
is not the spectrum of \(\mathcal E_s\).

The full trajectory-resolved record also defines classical information-theoretic
observables.  With
\begin{equation}
    Z_\xi=p_\xi,
    \qquad
    I_\xi=-\log p_\xi,
    \label{eq:app-record-self-information}
\end{equation}
the Shannon record entropy and the replicated sums are
\begin{align}
    H_{\rm rec}&=-\sum_\xi p_\xi\log p_\xi
    =\overline{I_\xi},\nonumber\\
    \mathcal Z_n&=\sum_\xi p_\xi^n,
    &H_n^{\rm rec}&=\frac{\log\mathcal Z_n}{1-n}.
    \label{eq:app-record-Renyi}
\end{align}
Here \(\overline{I_\xi}\) is exactly the Shannon self-information.  Naming it
a quenched free energy does not by itself supply a statistical-mechanics
theory.  The required bridge is
\begin{equation}
    -\left.\partial_n\log\mathcal Z_n\right|_{n=1}=H_{\rm rec},
    \label{eq:app-record-replica-derivative}
\end{equation}
together with an explicit replicated transfer normalization and boundary
conditions for \(\mathcal Z_n\).  Extracting a central charge or operator
content from finite-size corrections then requires an independent record
velocity, aspect-ratio convergence, and charge-resolved transfer sectors,
following, for example, Ref.~\cite{Zabalo2022}; it is not implied by the
state entanglement fit.


\section{Robustness protocol and numerical controls}
\label{app:robustness}

The numerical claims in the main text should be supported by a common set of
controls.  The following list defines the intended tests; no outcome is
assumed here.

\paragraph*{Initial states and convergence time.}
In addition to a generic pure Gaussian state and the maximally mixed state,
use onsite Fock product states
\begin{equation}
    G_{\rm prod}=\operatorname{diag}(2n_1-1,\ldots,2n_N-1),
    \qquad n_i\in\{0,1\}.
    \label{eq:app-product-initial-covariance}
\end{equation}
Vacuum, deterministic half-filled patterns, and random onsite occupations
test distinct charge sectors without importing initial spatial entanglement.
For each observable \(X\), define a tolerance-based convergence time
\begin{equation}
    \begin{aligned}
    T_\epsilon(X)=\min\bigl\{T:\ &
      \max_{T'\in[T,T+\Delta T]}\\[-0.25em]
      &|X(T')-X_{\rm late}|<\epsilon\bigr\}.
    \end{aligned}
    \label{eq:app-convergence-time}
\end{equation}
Its scaling with \(N_x,N_y\) must be measured rather than fixed by choosing
the final cycle proportional to system size.

\paragraph*{Locality and support.}
Repeat the analysis for several finite \(n_{\rm shell}\), and compare with
the dense frame only as a nonlocal reference.  A trend with support radius
must be separated from a trend with system size.  Compare the explicit
interface and support-terminated protocols as distinct microscopic
families.

\paragraph*{Schedule.}
Compare raster, randomized, and finite-color schedules at matched numbers of
local letters and at matched physical layer depth.  Record both the gate
count and the nonoverlapping circuit depth.  Universality under schedule
changes is stronger evidence than agreement after merely relabeling one
sweep as one unit of time.

\paragraph*{Imperfect feedback and ancilla noise.}
Let \(\epsilon_{\rm a}\) be the probability that the fresh ancilla is
prepared with the wrong occupation, and let \(p_{\rm g}\) and \(p_{\rm l}\)
be the success probabilities of gain and loss.  The ideal point is
\begin{equation}
    \epsilon_{\rm a}=0,
    \qquad p_{\rm g}=p_{\rm l}=1.
\end{equation}
Vary these independently, because a symmetric reduction can hide a
particle--hole-asymmetric instability.  A coherent fSWAP angle error
\(\theta=\pi/2+\delta\theta\) should be tested separately from classical
ancilla preparation error.

\paragraph*{Wall separation and aspect ratio.}
Vary both separations on the torus, not only \(N_y\).  Test the scaling
variable in Eq.~\eqref{eq:app-wall-width-condition} using an independently
measured \(\xi_\perp\), and compare against a domain-wall-off control at the
same dimensions and observation time.

\paragraph*{Tangent numerics.}
Report burn-in and observation windows, QR cadence, the number of propagated
vectors, singular tolerance, null counts, and censored trajectories.
Compare the direct chronological product at small size with QR accumulation.
At every fixed finite size, extend or extrapolate the QR observation time
before fitting any size dependence.  Randomize the initial tangent frame and
localize the projector onto the complete near-marginal subspace rather than
one basis-dependent QR vector.
For a half-strip, count \(N_=^\partial,N_{\ne}^\partial\) and verify
Eq.~\eqref{eq:app-record-resolved-cut-bound} from a pure product input.

\paragraph*{Fit and sampling stability.}
For entropy, charge fluctuations, and correlations, vary the fit window and
include the full trajectory bootstrap.  Report the number of independent
Born trajectories at every size.  Heatmaps are secondary localization
diagnostics and should not substitute for integrated wall weight, inverse
participation ratios, or transverse localization-length fits.

\appendixplaceholderfigure{fig:app-robustness}{Robustness and control matrix}{
Proposed panels: finite-support convergence; product, generic-Gaussian, and
maximally mixed initial states; raster versus finite-color schedules;
feedback and ancilla errors; wall-separation scaling; and null/censor
statistics for the tangent QR.  This is a placeholder layout only.  The final
figure should display all actual run parameters and should not combine data
from microscopically different domain-wall constructions without labels.}


\bibliographystyle{apsrev4-2}
\bibliography{../../LEGACY/topological_dynamics_introduction/references,references}

\end{document}
